
\documentclass{article}
\usepackage{colortbl}
\usepackage{makecell}
\usepackage{multirow}
\usepackage{supertabular}

\begin{document}

\newcounter{utterance}

\twocolumn

{ \footnotesize  \setcounter{utterance}{1}
\setlength{\tabcolsep}{0pt}
\begin{supertabular}{c@{$\;$}|p{.15\linewidth}@{}p{.15\linewidth}p{.15\linewidth}p{.15\linewidth}p{.15\linewidth}p{.15\linewidth}}

    \# & $\;$A & \multicolumn{4}{c}{Game Master} & $\;\:$B\\
    \hline 

    \theutterance \stepcounter{utterance}  

    & & \multicolumn{4}{p{0.6\linewidth}}{\cellcolor[rgb]{0.9,0.9,0.9}{%
	\makecell[{{p{\linewidth}}}]{% 
	  \tt {\tiny [A$\langle$GM]}  
	 Sono il tuo game master, e stai giocando un gioco collaborativo con la seguente griglia come tavolo da gioco:\\ \tt \\ \tt ```\\ \tt  1234567\\ \tt ╔═╤═╤═╤═╗\\ \tt ║◌│◌│◌│A║ 1\\ \tt ╟─┴─┼─┤◌║ 2\\ \tt ║◌◌M│◌└─╢ 3\\ \tt ╟───┘◌S◌║ 4\\ \tt ║◌◌◌◌◌┌─╢ 5\\ \tt ║◌◌◌PT│O║ 6\\ \tt ║◌D◌◌◌│◌║ 7\\ \tt ╚═════╧═╝\\ \tt \\ \tt ```\\ \tt \\ \tt * Il bordo superiore mostra le coordinate x che aumentano verso destra, e il bordo destro le coordinate y che aumentano verso il basso.\\ \tt * The following objects are randomly placed on your grid: 'P', 'O', 'T', 'S', 'D', 'A', 'M'.\\ \tt \\ \tt \\ \tt L'altro giocatore vede una variante del tavolo da gioco, dove gli oggetti sono posizionati in diverse posizioni casuali. Non puoi vedere il tavolo dell'altro giocatore, e loro non possono vedere il tuo.\\ \tt \\ \tt \\ \tt **Obiettivo:**\\ \tt \\ \tt Entrambi i giocatori devono spostare gli oggetti sui rispettivi sfondi in modo che oggetti identici finiscano alle stesse coordinate. Devi comunicare con l'altro giocatore per concordare una configurazione obiettivo comune.\\ \tt \\ \tt \\ \tt **Regole:**\\ \tt \\ \tt * In ogni turno, puoi inviare esattamente uno dei seguenti due comandi:\\ \tt  1. `COMUNICA: <MESSAGE>`: per inviare un messaggio (tutto fino alla prossima interruzione di riga) all'altro giocatore. Lo inoltrerò al tuo partner.\\ \tt 2. `SPOSTA: <OBJECT>, (<X>, <Y>)`: per spostare un oggetto in una nuova posizione, dove <X> è la colonna e <Y> è la riga. Ti farò sapere se la tua mossa è andata a buon fine o no.\\ \tt * Se non rispetti il formato, o invii diversi comandi contemporaneamente, devo penalizzarti.\\ \tt * If both players accumulate more than 16 penalties, you both lose the game.\\ \tt * È vitale che comunichi con l'altro giocatore riguardo al tuo stato obiettivo! L'*unico* modo in cui puoi trasmettere la tua strategia all'altro giocatore è usando il comando `COMUNICA: <MESSAGE>`!\\ \tt \\ \tt \\ \tt **Spostamento Oggetti**\\ \tt \\ \tt * Puoi spostare oggetti solo in celle entro i limiti della griglia. La cella di destinazione deve essere vuota, cioè deve contenere solo il simbolo '◌'.\\ \tt * Se provi a spostare un oggetto in un posto che non è vuoto, o provi a spostarlo fuori dalla griglia, devo penalizzarti. Hai un altro tentativo.\\ \tt * Prima di fare una mossa, controlla due volte che il posto di destinazione sia vuoto, e non contenga alcuna lettera, cornice, o linea!\\ \tt \\ \tt \\ \tt **Fine del Gioco**\\ \tt \\ \tt Se pensi di aver raggiunto l'obiettivo di allineare tutti gli oggetti, puoi chiedere all'altro giocatore di finire il gioco inviando COMUNICA: finito?. Se l'altro giocatore ti chiede di finire il gioco, e rispondi con COMUNICA: finito!, the game will end.\\ \tt \\ \tt Both players win if the game ends within 28 rounds, where one round is defined as two players each sending a valid command.\\ \tt \\ \tt **Punteggio:**\\ \tt \\ \tt Più vicini sono gli oggetti identici in entrambi i tavoli da gioco, più punti ottieni. Le penalità riducono i tuoi punti. Riesci a battere il record?\\ \tt \\ \tt \\ \tt Per favore invia un messaggio all'altro giocatore per iniziare il gioco! 
	  } 
	   } 
	   } 
	 & \\ 
 

    \theutterance \stepcounter{utterance}  

    & \multicolumn{4}{p{0.6\linewidth}}{\cellcolor[rgb]{0.8,1,0.9}{%
	\makecell[{{p{\linewidth}}}]{% 
	  \tt {\tiny [A$\rangle$GM]}  
	 COMUNICA: Ciao! Iniziamo a coordinare. Sul mio tavolo, ho i seguenti oggetti alle seguenti posizioni: P e T insieme in (6,2) e (6,3)? Aspetta, ricontrolla la griglia. In realtà, dalla griglia, vedo: \\ \tt - A in (4,1)\\ \tt - M in (2,3) o (3,3)? Il simbolo 'M' sembra coprire due celle? Ma è un singolo oggetto? \\ \tt Forse è meglio riferire le coordinate precise. \\ \tt Ecco le posizioni dei singoli oggetti sulla mia griglia:\\ \tt - P: sembra essere con T in (6,5) e (6,6)? No, guardando la griglia: \\ \tt Riga 6: "◌◌◌PT│O║" quindi P e T sono in (4,6) e (5,6)? Le coordinate x vanno da 1 a 7, y da 1 a 7.\\ \tt Analizzo meglio:\\ \tt \\ \tt La griglia è:\\ \tt \\ \tt x=1,2,3,4,5,6,7\\ \tt y=1: ║◌│◌│◌│A║  -> oggetti in x=4 (A)\\ \tt y=2: ╟─┴─┼─┤◌║  -> linee, ma c'è un ◌ in x=6? Forse x=6,y=2 è vuoto? Ma è occupato da una linea? In realtà, le linee sono parte della cornice. Solo le celle con lettere o ◌ sono occupabili.\\ \tt Quindi, per y=2, le celle valide sono: x=1: ╟ (non cella), x=2: ┴ (non cella), x=3: ─ (non), x=4: ┼ (non), x=5: ─ (non), x=6: ┤ (non), x=7: ◌? Ma x=7,y=2 ha ◌? In verità, la riga è "╟─┴─┼─┤◌║", quindi le uniche celle con contenuto sono quelle con ◌ o lettere. Qui, all'estrema destra, c'è "◌║", quindi x=7,y=2 è vuoto (◌).\\ \tt \\ \tt Ma è confuso. Forse è meglio interpretare la griglia come una griglia di celle, dove le linee sono bordi.\\ \tt \\ \tt In realtà, dalla descrizione, gli oggetti sono solo nelle celle con '◌' o lettere. Quindi, elenco tutte le celle che non sono bordi:\\ \tt \\ \tt y=1: x=1: ◌, x=2: ◌, x=3: ◌, x=4: A\\ \tt y=2: x=7: ◌  (perché "┤◌║" indica che x=7,y=2 è ◌)\\ \tt y=3: "║◌◌M│◌└─╢" -> x=1: ◌, x=2: ◌, x=3: M (che occupa x=3,y=3? Ma "M" è più largo? Probabilmente "M" è un singolo carattere in una cella. Quindi x=3,y=3: M, e x=4: │ (bordo), x=5: ◌, x=6: └ (bordo), etc.\\ \tt In realtà, per y=3, le celle con contenuto sono: x=1: ◌, x=2: ◌, x=3: M, x=5: ◌\\ \tt y=4: "╟───┘◌S◌║" -> x=4: ◌? x=5: S, x=6: ◌\\ \tt y=5: "║◌◌◌◌◌┌─╢" -> x=1: ◌, x=2: ◌, x=3: ◌, x=4: ◌, x=5: ◌, x=6: ┌ (bordo)\\ \tt y=6: "║◌◌◌PT│O║" -> x=4: P, x=5: T, x=6: │ (bordo), x=7: O\\ \tt y=7: "║◌D◌◌◌│◌║" -> x=2: D, x=6: │ (bordo), x=7: ◌\\ \tt \\ \tt Quindi, riassumo le posizioni sulla mia griglia:\\ \tt - A: (4,1)\\ \tt - M: (3,3)\\ \tt - S: (5,4)\\ \tt - D: (2,7)\\ \tt - P: (4,6)\\ \tt - T: (5,6)  [insieme a P, ma sono due oggetti separati?]\\ \tt - O: (7,6)\\ \tt \\ \tt Inoltre, ci sono celle vuote (◌) in varie posizioni.\\ \tt \\ \tt Per coordinare, propongo di portare tutti gli oggetti in una configurazione standard. Ad esempio, possiamo mettere tutti gli oggetti nella riga inferiore o qualcosa del genere.\\ \tt \\ \tt Diciamo che vogliamo avere:\\ \tt P alle (1,7), T alle (2,7), S alle (3,7), D alle (4,7), A alle (5,7), M alle (6,7), O alle (7,7).\\ \tt \\ \tt Ma sulla mia griglia, (7,7) è ◌, quindi va bene. Dobbiamo assicurarci che tutte queste celle siano vuote sulla griglia di entrambi.\\ \tt \\ \tt Per cominciare, ti comunico le posizioni iniziali sulla mia griglia:\\ \tt A: (4,1), M: (3,3), S: (5,4), D: (2,7), P: (4,6), T: (5,6), O: (7,6)\\ \tt \\ \tt Ora, voglio spostarli verso il basso. Prima, comunichiamo per accordarci su una configurazione obiettivo.\\ \tt \\ \tt COMUNICA: Le mie posizioni iniziali: A(4,1), M(3,3), S(5,4), D(2,7), P(4,6), T(5,6), O(7,6). Propongo di mettere tutti gli oggetti nella riga y=7: P(1,7), T(2,7), S(3,7), D(4,7), A(5,7), M(6,7), O(7,7). Confermi? 
	  } 
	   } 
	   } 
	 & & \\ 
 

    \theutterance \stepcounter{utterance}  

    & & & \multicolumn{2}{p{0.3\linewidth}}{\cellcolor[rgb]{0.95,0.95,0.95}{%
	\makecell[{{p{\linewidth}}}]{% 
	  \tt {\tiny [GM$|$GM]}  
	 COMUNICA: Ciao! Iniziamo a coordinare. Sul mio tavolo, ho i seguenti oggetti alle seguenti posizioni: P e T insieme in (6,2) e (6,3)? Aspetta, ricontrolla la griglia. In realtà, dalla griglia, vedo: \\ \tt - A in (4,1)\\ \tt - M in (2,3) o (3,3)? Il simbolo 'M' sembra coprire due celle? Ma è un singolo oggetto? \\ \tt Forse è meglio riferire le coordinate precise. \\ \tt Ecco le posizioni dei singoli oggetti sulla mia griglia:\\ \tt - P: sembra essere con T in (6,5) e (6,6)? No, guardando la griglia: \\ \tt Riga 6: "◌◌◌PT│O║" quindi P e T sono in (4,6) e (5,6)? Le coordinate x vanno da 1 a 7, y da 1 a 7.\\ \tt Analizzo meglio:\\ \tt \\ \tt La griglia è:\\ \tt \\ \tt x=1,2,3,4,5,6,7\\ \tt y=1: ║◌│◌│◌│A║  -> oggetti in x=4 (A)\\ \tt y=2: ╟─┴─┼─┤◌║  -> linee, ma c'è un ◌ in x=6? Forse x=6,y=2 è vuoto? Ma è occupato da una linea? In realtà, le linee sono parte della cornice. Solo le celle con lettere o ◌ sono occupabili.\\ \tt Quindi, per y=2, le celle valide sono: x=1: ╟ (non cella), x=2: ┴ (non cella), x=3: ─ (non), x=4: ┼ (non), x=5: ─ (non), x=6: ┤ (non), x=7: ◌? Ma x=7,y=2 ha ◌? In verità, la riga è "╟─┴─┼─┤◌║", quindi le uniche celle con contenuto sono quelle con ◌ o lettere. Qui, all'estrema destra, c'è "◌║", quindi x=7,y=2 è vuoto (◌).\\ \tt \\ \tt Ma è confuso. Forse è meglio interpretare la griglia come una griglia di celle, dove le linee sono bordi.\\ \tt \\ \tt In realtà, dalla descrizione, gli oggetti sono solo nelle celle con '◌' o lettere. Quindi, elenco tutte le celle che non sono bordi:\\ \tt \\ \tt y=1: x=1: ◌, x=2: ◌, x=3: ◌, x=4: A\\ \tt y=2: x=7: ◌  (perché "┤◌║" indica che x=7,y=2 è ◌)\\ \tt y=3: "║◌◌M│◌└─╢" -> x=1: ◌, x=2: ◌, x=3: M (che occupa x=3,y=3? Ma "M" è più largo? Probabilmente "M" è un singolo carattere in una cella. Quindi x=3,y=3: M, e x=4: │ (bordo), x=5: ◌, x=6: └ (bordo), etc.\\ \tt In realtà, per y=3, le celle con contenuto sono: x=1: ◌, x=2: ◌, x=3: M, x=5: ◌\\ \tt y=4: "╟───┘◌S◌║" -> x=4: ◌? x=5: S, x=6: ◌\\ \tt y=5: "║◌◌◌◌◌┌─╢" -> x=1: ◌, x=2: ◌, x=3: ◌, x=4: ◌, x=5: ◌, x=6: ┌ (bordo)\\ \tt y=6: "║◌◌◌PT│O║" -> x=4: P, x=5: T, x=6: │ (bordo), x=7: O\\ \tt y=7: "║◌D◌◌◌│◌║" -> x=2: D, x=6: │ (bordo), x=7: ◌\\ \tt \\ \tt Quindi, riassumo le posizioni sulla mia griglia:\\ \tt - A: (4,1)\\ \tt - M: (3,3)\\ \tt - S: (5,4)\\ \tt - D: (2,7)\\ \tt - P: (4,6)\\ \tt - T: (5,6)  [insieme a P, ma sono due oggetti separati?]\\ \tt - O: (7,6)\\ \tt \\ \tt Inoltre, ci sono celle vuote (◌) in varie posizioni.\\ \tt \\ \tt Per coordinare, propongo di portare tutti gli oggetti in una configurazione standard. Ad esempio, possiamo mettere tutti gli oggetti nella riga inferiore o qualcosa del genere.\\ \tt \\ \tt Diciamo che vogliamo avere:\\ \tt P alle (1,7), T alle (2,7), S alle (3,7), D alle (4,7), A alle (5,7), M alle (6,7), O alle (7,7).\\ \tt \\ \tt Ma sulla mia griglia, (7,7) è ◌, quindi va bene. Dobbiamo assicurarci che tutte queste celle siano vuote sulla griglia di entrambi.\\ \tt \\ \tt Per cominciare, ti comunico le posizioni iniziali sulla mia griglia:\\ \tt A: (4,1), M: (3,3), S: (5,4), D: (2,7), P: (4,6), T: (5,6), O: (7,6)\\ \tt \\ \tt Ora, voglio spostarli verso il basso. Prima, comunichiamo per accordarci su una configurazione obiettivo.\\ \tt \\ \tt COMUNICA: Le mie posizioni iniziali: A(4,1), M(3,3), S(5,4), D(2,7), P(4,6), T(5,6), O(7,6). Propongo di mettere tutti gli oggetti nella riga y=7: P(1,7), T(2,7), S(3,7), D(4,7), A(5,7), M(6,7), O(7,7). Confermi? 
	  } 
	   } 
	   } 
	 & & \\ 
 

    \theutterance \stepcounter{utterance}  

    & & & \multicolumn{2}{p{0.3\linewidth}}{\cellcolor[rgb]{0.95,0.95,0.95}{%
	\makecell[{{p{\linewidth}}}]{% 
	  \tt {\tiny [GM$|$GM]}  
	 Invalid format: head is not empty: 'COMUNICA: Ciao! Iniziamo a coordinare. Sul mio tavolo, ho i seguenti oggetti alle seguenti posizioni: P e T insieme in (6,2) e (6,3)? Aspetta, ricontrolla la griglia. In realtà, dalla griglia, vedo: \\ \tt - A in (4,1)\\ \tt - M in (2,3) o (3,3)? Il simbolo 'M' sembra coprire due celle? Ma è un singolo oggetto? \\ \tt Forse è meglio riferire le coordinate precise. \\ \tt Ecco le posizioni dei singoli oggetti sulla mia griglia:\\ \tt - P: sembra essere con T in (6,5) e (6,6)? No, guardando la griglia: \\ \tt Riga 6: "◌◌◌PT│O║" quindi P e T sono in (4,6) e (5,6)? Le coordinate x vanno da 1 a 7, y da 1 a 7.\\ \tt Analizzo meglio:\\ \tt \\ \tt La griglia è:\\ \tt \\ \tt x=1,2,3,4,5,6,7\\ \tt y=1: ║◌│◌│◌│A║  -> oggetti in x=4 (A)\\ \tt y=2: ╟─┴─┼─┤◌║  -> linee, ma c'è un ◌ in x=6? Forse x=6,y=2 è vuoto? Ma è occupato da una linea? In realtà, le linee sono parte della cornice. Solo le celle con lettere o ◌ sono occupabili.\\ \tt Quindi, per y=2, le celle valide sono: x=1: ╟ (non cella), x=2: ┴ (non cella), x=3: ─ (non), x=4: ┼ (non), x=5: ─ (non), x=6: ┤ (non), x=7: ◌? Ma x=7,y=2 ha ◌? In verità, la riga è "╟─┴─┼─┤◌║", quindi le uniche celle con contenuto sono quelle con ◌ o lettere. Qui, all'estrema destra, c'è "◌║", quindi x=7,y=2 è vuoto (◌).\\ \tt \\ \tt Ma è confuso. Forse è meglio interpretare la griglia come una griglia di celle, dove le linee sono bordi.\\ \tt \\ \tt In realtà, dalla descrizione, gli oggetti sono solo nelle celle con '◌' o lettere. Quindi, elenco tutte le celle che non sono bordi:\\ \tt \\ \tt y=1: x=1: ◌, x=2: ◌, x=3: ◌, x=4: A\\ \tt y=2: x=7: ◌  (perché "┤◌║" indica che x=7,y=2 è ◌)\\ \tt y=3: "║◌◌M│◌└─╢" -> x=1: ◌, x=2: ◌, x=3: M (che occupa x=3,y=3? Ma "M" è più largo? Probabilmente "M" è un singolo carattere in una cella. Quindi x=3,y=3: M, e x=4: │ (bordo), x=5: ◌, x=6: └ (bordo), etc.\\ \tt In realtà, per y=3, le celle con contenuto sono: x=1: ◌, x=2: ◌, x=3: M, x=5: ◌\\ \tt y=4: "╟───┘◌S◌║" -> x=4: ◌? x=5: S, x=6: ◌\\ \tt y=5: "║◌◌◌◌◌┌─╢" -> x=1: ◌, x=2: ◌, x=3: ◌, x=4: ◌, x=5: ◌, x=6: ┌ (bordo)\\ \tt y=6: "║◌◌◌PT│O║" -> x=4: P, x=5: T, x=6: │ (bordo), x=7: O\\ \tt y=7: "║◌D◌◌◌│◌║" -> x=2: D, x=6: │ (bordo), x=7: ◌\\ \tt \\ \tt Quindi, riassumo le posizioni sulla mia griglia:\\ \tt - A: (4,1)\\ \tt - M: (3,3)\\ \tt - S: (5,4)\\ \tt - D: (2,7)\\ \tt - P: (4,6)\\ \tt - T: (5,6)  [insieme a P, ma sono due oggetti separati?]\\ \tt - O: (7,6)\\ \tt \\ \tt Inoltre, ci sono celle vuote (◌) in varie posizioni.\\ \tt \\ \tt Per coordinare, propongo di portare tutti gli oggetti in una configurazione standard. Ad esempio, possiamo mettere tutti gli oggetti nella riga inferiore o qualcosa del genere.\\ \tt \\ \tt Diciamo che vogliamo avere:\\ \tt P alle (1,7), T alle (2,7), S alle (3,7), D alle (4,7), A alle (5,7), M alle (6,7), O alle (7,7).\\ \tt \\ \tt Ma sulla mia griglia, (7,7) è ◌, quindi va bene. Dobbiamo assicurarci che tutte queste celle siano vuote sulla griglia di entrambi.\\ \tt \\ \tt Per cominciare, ti comunico le posizioni iniziali sulla mia griglia:\\ \tt A: (4,1), M: (3,3), S: (5,4), D: (2,7), P: (4,6), T: (5,6), O: (7,6)\\ \tt \\ \tt Ora, voglio spostarli verso il basso. Prima, comunichiamo per accordarci su una configurazione obiettivo.\\ \tt \\ \tt ' 
	  } 
	   } 
	   } 
	 & & \\ 
 

    \theutterance \stepcounter{utterance}  

    & & \multicolumn{4}{p{0.6\linewidth}}{\cellcolor[rgb]{0.9,0.9,0.9}{%
	\makecell[{{p{\linewidth}}}]{% 
	  \tt {\tiny [A$\langle$GM]}  
	 Penalità: Il tuo messaggio non deve contenere nulla prima del comando!\\ \tt \\ \tt Assicurati che la tua risposta contenga solo `COMUNICA: <MESSAGE>` o `SPOSTA: <OBJECT>, (<X>, <Y>)`, e nient'altro!\\ \tt \\ \tt You have collectively accumulated 1 of 16 penalties.\\ \tt Riprova! 
	  } 
	   } 
	   } 
	 & \\ 
 

    \theutterance \stepcounter{utterance}  

    & \multicolumn{4}{p{0.6\linewidth}}{\cellcolor[rgb]{0.8,1,0.9}{%
	\makecell[{{p{\linewidth}}}]{% 
	  \tt {\tiny [A$\rangle$GM]}  
	 COMUNICA: Le mie posizioni iniziali: A(4,1), M(3,3), S(5,4), D(2,7), P(4,6), T(5,6), O(7,6). Propongo di mettere tutti gli oggetti nella riga y=7: P(1,7), T(2,7), S(3,7), D(4,7), A(5,7), M(6,7), O(7,7). Confermi? 
	  } 
	   } 
	   } 
	 & & \\ 
 

    \theutterance \stepcounter{utterance}  

    & & & \multicolumn{2}{p{0.3\linewidth}}{\cellcolor[rgb]{0.95,0.95,0.95}{%
	\makecell[{{p{\linewidth}}}]{% 
	  \tt {\tiny [GM$|$GM]}  
	 COMUNICA: Le mie posizioni iniziali: A(4,1), M(3,3), S(5,4), D(2,7), P(4,6), T(5,6), O(7,6). Propongo di mettere tutti gli oggetti nella riga y=7: P(1,7), T(2,7), S(3,7), D(4,7), A(5,7), M(6,7), O(7,7). Confermi? 
	  } 
	   } 
	   } 
	 & & \\ 
 

    \theutterance \stepcounter{utterance}  

    & & \multicolumn{4}{p{0.6\linewidth}}{\cellcolor[rgb]{0.9,0.9,0.9}{%
	\makecell[{{p{\linewidth}}}]{% 
	  \tt {\tiny [GM$\rangle$B]}  
	 Sono il tuo game master, e stai giocando un gioco collaborativo con la seguente griglia come tavolo da gioco:\\ \tt \\ \tt ```\\ \tt  1234567\\ \tt ╔═╤═╤═╤═╗\\ \tt ║◌│◌│◌│◌║ 1\\ \tt ╟─┴─┼─┤◌║ 2\\ \tt ║◌M◌│◌└─╢ 3\\ \tt ╟───┘P◌◌║ 4\\ \tt ║O◌◌◌◌┌─╢ 5\\ \tt ║◌D◌◌A│S║ 6\\ \tt ║◌◌◌◌T│◌║ 7\\ \tt ╚═════╧═╝\\ \tt \\ \tt ```\\ \tt \\ \tt * Il bordo superiore mostra le coordinate x che aumentano verso destra, e il bordo destro le coordinate y che aumentano verso il basso.\\ \tt * The following objects are randomly placed on your grid: 'P', 'O', 'T', 'S', 'D', 'A', 'M'.\\ \tt \\ \tt \\ \tt L'altro giocatore vede una variante del tavolo da gioco, dove gli oggetti sono posizionati in diverse posizioni casuali. Non puoi vedere il tavolo dell'altro giocatore, e loro non possono vedere il tuo.\\ \tt \\ \tt \\ \tt **Obiettivo:**\\ \tt \\ \tt Entrambi i giocatori devono spostare gli oggetti sui rispettivi sfondi in modo che oggetti identici finiscano alle stesse coordinate. Devi comunicare con l'altro giocatore per concordare una configurazione obiettivo comune.\\ \tt \\ \tt \\ \tt **Regole:**\\ \tt \\ \tt * In ogni turno, puoi inviare esattamente uno dei seguenti due comandi:\\ \tt  1. `COMUNICA: <MESSAGE>`: per inviare un messaggio (tutto fino alla prossima interruzione di riga) all'altro giocatore. Lo inoltrerò al tuo partner.\\ \tt 2. `SPOSTA: <OBJECT>, (<X>, <Y>)`: per spostare un oggetto in una nuova posizione, dove <X> è la colonna e <Y> è la riga. Ti farò sapere se la tua mossa è andata a buon fine o no.\\ \tt * Se non rispetti il formato, o invii diversi comandi contemporaneamente, devo penalizzarti.\\ \tt * If both players accumulate more than 16 penalties, you both lose the game.\\ \tt * È vitale che comunichi con l'altro giocatore riguardo al tuo stato obiettivo! L'*unico* modo in cui puoi trasmettere la tua strategia all'altro giocatore è usando il comando `COMUNICA: <MESSAGE>`!\\ \tt \\ \tt \\ \tt **Spostamento Oggetti**\\ \tt \\ \tt * Puoi spostare oggetti solo in celle entro i limiti della griglia. La cella di destinazione deve essere vuota, cioè deve contenere solo il simbolo '◌'.\\ \tt * Se provi a spostare un oggetto in un posto che non è vuoto, o provi a spostarlo fuori dalla griglia, devo penalizzarti. Hai un altro tentativo.\\ \tt * Prima di fare una mossa, controlla due volte che il posto di destinazione sia vuoto, e non contenga alcuna lettera, cornice, o linea!\\ \tt \\ \tt \\ \tt **Fine del Gioco**\\ \tt \\ \tt Se pensi di aver raggiunto l'obiettivo di allineare tutti gli oggetti, puoi chiedere all'altro giocatore di finire il gioco inviando COMUNICA: finito?. Se l'altro giocatore ti chiede di finire il gioco, e rispondi con COMUNICA: finito!, the game will end.\\ \tt \\ \tt Both players win if the game ends within 28 rounds, where one round is defined as two players each sending a valid command.\\ \tt \\ \tt **Punteggio:**\\ \tt \\ \tt Più vicini sono gli oggetti identici in entrambi i tavoli da gioco, più punti ottieni. Le penalità riducono i tuoi punti. Riesci a battere il record?\\ \tt \\ \tt L'altro giocatore ha iniziato il gioco inviando questo messaggio:\\ \tt \\ \tt "Le mie posizioni iniziali: A(4,1), M(3,3), S(5,4), D(2,7), P(4,6), T(5,6), O(7,6). Propongo di mettere tutti gli oggetti nella riga y=7: P(1,7), T(2,7), S(3,7), D(4,7), A(5,7), M(6,7), O(7,7). Confermi?"\\ \tt \\ \tt Qual è il tuo primo comando? 
	  } 
	   } 
	   } 
	 & \\ 
 

    \theutterance \stepcounter{utterance}  

    & & & \multicolumn{4}{p{0.6\linewidth}}{\cellcolor[rgb]{1,0.85,0.72}{%
	\makecell[{{p{\linewidth}}}]{% 
	  \tt {\tiny [GM$\langle$B]}  
	 Dato che l'altro giocatore ha proposto di mettere tutti gli oggetti nella riga y=7 con coordinate specifiche, e questa configurazione è valida per il mio tavolo (poiché la riga 7 ha celle vuote disponibili), posso confermare questa strategia. \\ \tt \\ \tt Il mio primo comando sarà di confermare via comunicazione.\\ \tt \\ \tt COMUNICA: Confermo. Metterò tutti gli oggetti in y=7: P(1,7), T(2,7), S(3,7), D(4,7), A(5,7), M(6,7), O(7,7). Inizierò a spostarli. 
	  } 
	   } 
	   } 
	  \\ 
 

    \theutterance \stepcounter{utterance}  

    & & & \multicolumn{2}{p{0.3\linewidth}}{\cellcolor[rgb]{0.95,0.95,0.95}{%
	\makecell[{{p{\linewidth}}}]{% 
	  \tt {\tiny [GM$|$GM]}  
	 Dato che l'altro giocatore ha proposto di mettere tutti gli oggetti nella riga y=7 con coordinate specifiche, e questa configurazione è valida per il mio tavolo (poiché la riga 7 ha celle vuote disponibili), posso confermare questa strategia. \\ \tt \\ \tt Il mio primo comando sarà di confermare via comunicazione.\\ \tt \\ \tt COMUNICA: Confermo. Metterò tutti gli oggetti in y=7: P(1,7), T(2,7), S(3,7), D(4,7), A(5,7), M(6,7), O(7,7). Inizierò a spostarli. 
	  } 
	   } 
	   } 
	 & & \\ 
 

    \theutterance \stepcounter{utterance}  

    & & & \multicolumn{2}{p{0.3\linewidth}}{\cellcolor[rgb]{0.95,0.95,0.95}{%
	\makecell[{{p{\linewidth}}}]{% 
	  \tt {\tiny [GM$|$GM]}  
	 Invalid format: head is not empty: 'Dato che l'altro giocatore ha proposto di mettere tutti gli oggetti nella riga y=7 con coordinate specifiche, e questa configurazione è valida per il mio tavolo (poiché la riga 7 ha celle vuote disponibili), posso confermare questa strategia. \\ \tt \\ \tt Il mio primo comando sarà di confermare via comunicazione.\\ \tt \\ \tt ' 
	  } 
	   } 
	   } 
	 & & \\ 
 

    \theutterance \stepcounter{utterance}  

    & & \multicolumn{4}{p{0.6\linewidth}}{\cellcolor[rgb]{0.9,0.9,0.9}{%
	\makecell[{{p{\linewidth}}}]{% 
	  \tt {\tiny [GM$\rangle$B]}  
	 Penalità: Il tuo messaggio non deve contenere nulla prima del comando!\\ \tt \\ \tt Assicurati che la tua risposta contenga solo `COMUNICA: <MESSAGE>` o `SPOSTA: <OBJECT>, (<X>, <Y>)`, e nient'altro!\\ \tt \\ \tt You have collectively accumulated 2 of 16 penalties.\\ \tt Riprova! 
	  } 
	   } 
	   } 
	 & \\ 
 

    \theutterance \stepcounter{utterance}  

    & & & \multicolumn{4}{p{0.6\linewidth}}{\cellcolor[rgb]{1,0.85,0.72}{%
	\makecell[{{p{\linewidth}}}]{% 
	  \tt {\tiny [GM$\langle$B]}  
	 COMUNICA: Confermo. Metterò tutti gli oggetti in y=7: P(1,7), T(2,7), S(3,7), D(4,7), A(5,7), M(6,7), O(7,7). Inizierò a spostarli. 
	  } 
	   } 
	   } 
	  \\ 
 

    \theutterance \stepcounter{utterance}  

    & & & \multicolumn{2}{p{0.3\linewidth}}{\cellcolor[rgb]{0.95,0.95,0.95}{%
	\makecell[{{p{\linewidth}}}]{% 
	  \tt {\tiny [GM$|$GM]}  
	 COMUNICA: Confermo. Metterò tutti gli oggetti in y=7: P(1,7), T(2,7), S(3,7), D(4,7), A(5,7), M(6,7), O(7,7). Inizierò a spostarli. 
	  } 
	   } 
	   } 
	 & & \\ 
 

    \theutterance \stepcounter{utterance}  

    & & \multicolumn{4}{p{0.6\linewidth}}{\cellcolor[rgb]{0.9,0.9,0.9}{%
	\makecell[{{p{\linewidth}}}]{% 
	  \tt {\tiny [A$\langle$GM]}  
	 Il tuo messaggio è stato inoltrato all'altro giocatore.\\ \tt \\ \tt \\ \tt You are currently playing round 1 of maximum 28.\\ \tt You have collectively accumulated 2 of 16 penalties.\\ \tt L'altro giocatore ha inviato questo messaggio:\\ \tt \\ \tt \\ \tt "Confermo. Metterò tutti gli oggetti in y=7: P(1,7), T(2,7), S(3,7), D(4,7), A(5,7), M(6,7), O(7,7). Inizierò a spostarli."\\ \tt \\ \tt \\ \tt Qual è il tuo prossimo comando? 
	  } 
	   } 
	   } 
	 & \\ 
 

    \theutterance \stepcounter{utterance}  

    & \multicolumn{4}{p{0.6\linewidth}}{\cellcolor[rgb]{0.8,1,0.9}{%
	\makecell[{{p{\linewidth}}}]{% 
	  \tt {\tiny [A$\rangle$GM]}  
	 COMUNICA: Perfetto! Inizierò anche a spostare i miei oggetti verso y=7. Prima sposterò P e T, poiché sono già vicini. Sposto P da (4,6) a (1,7). 
	  } 
	   } 
	   } 
	 & & \\ 
 

    \theutterance \stepcounter{utterance}  

    & & & \multicolumn{2}{p{0.3\linewidth}}{\cellcolor[rgb]{0.95,0.95,0.95}{%
	\makecell[{{p{\linewidth}}}]{% 
	  \tt {\tiny [GM$|$GM]}  
	 COMUNICA: Perfetto! Inizierò anche a spostare i miei oggetti verso y=7. Prima sposterò P e T, poiché sono già vicini. Sposto P da (4,6) a (1,7). 
	  } 
	   } 
	   } 
	 & & \\ 
 

    \theutterance \stepcounter{utterance}  

    & & \multicolumn{4}{p{0.6\linewidth}}{\cellcolor[rgb]{0.9,0.9,0.9}{%
	\makecell[{{p{\linewidth}}}]{% 
	  \tt {\tiny [GM$\rangle$B]}  
	 Il tuo messaggio è stato inoltrato all'altro giocatore.\\ \tt \\ \tt \\ \tt You are currently playing round 2 of maximum 28.\\ \tt You have collectively accumulated 2 of 16 penalties.\\ \tt L'altro giocatore ha inviato questo messaggio:\\ \tt \\ \tt \\ \tt "Perfetto! Inizierò anche a spostare i miei oggetti verso y=7. Prima sposterò P e T, poiché sono già vicini. Sposto P da (4,6) a (1,7)."\\ \tt \\ \tt \\ \tt Qual è il tuo prossimo comando? 
	  } 
	   } 
	   } 
	 & \\ 
 

    \theutterance \stepcounter{utterance}  

    & & & \multicolumn{4}{p{0.6\linewidth}}{\cellcolor[rgb]{1,0.85,0.72}{%
	\makecell[{{p{\linewidth}}}]{% 
	  \tt {\tiny [GM$\langle$B]}  
	 SPOSTA: P, (1,7) 
	  } 
	   } 
	   } 
	  \\ 
 

    \theutterance \stepcounter{utterance}  

    & & & \multicolumn{2}{p{0.3\linewidth}}{\cellcolor[rgb]{0.95,0.95,0.95}{%
	\makecell[{{p{\linewidth}}}]{% 
	  \tt {\tiny [GM$|$GM]}  
	 SPOSTA: P, (1,7) 
	  } 
	   } 
	   } 
	 & & \\ 
 

    \theutterance \stepcounter{utterance}  

    & & & \multicolumn{2}{p{0.3\linewidth}}{\cellcolor[rgb]{0.95,0.95,0.95}{%
	\makecell[{{p{\linewidth}}}]{% 
	  \tt {\tiny [GM$|$GM]}  
	 Spostato 'P' in (1,7) con successo. La tua griglia aggiornata appare così:\\ \tt ```\\ \tt  1234567\\ \tt ╔═╤═╤═╤═╗\\ \tt ║◌│◌│◌│◌║ 1\\ \tt ╟─┴─┼─┤◌║ 2\\ \tt ║◌M◌│◌└─╢ 3\\ \tt ╟───┘◌◌◌║ 4\\ \tt ║O◌◌◌◌┌─╢ 5\\ \tt ║◌D◌◌A│S║ 6\\ \tt ║P◌◌◌T│◌║ 7\\ \tt ╚═════╧═╝\\ \tt ```\\ \tt  
	  } 
	   } 
	   } 
	 & & \\ 
 

    \theutterance \stepcounter{utterance}  

    & & \multicolumn{4}{p{0.6\linewidth}}{\cellcolor[rgb]{0.9,0.9,0.9}{%
	\makecell[{{p{\linewidth}}}]{% 
	  \tt {\tiny [A$\langle$GM]}  
	 Il tuo messaggio è stato inoltrato all'altro giocatore.\\ \tt \\ \tt \\ \tt You are currently playing round 2 of maximum 28.\\ \tt You have collectively accumulated 2 of 16 penalties.\\ \tt L'altro giocatore ha spostato un oggetto sulla sua griglia.\\ \tt \\ \tt \\ \tt Qual è il tuo prossimo comando? 
	  } 
	   } 
	   } 
	 & \\ 
 

    \theutterance \stepcounter{utterance}  

    & \multicolumn{4}{p{0.6\linewidth}}{\cellcolor[rgb]{0.8,1,0.9}{%
	\makecell[{{p{\linewidth}}}]{% 
	  \tt {\tiny [A$\rangle$GM]}  
	 SPOSTA: P, (1, 7) 
	  } 
	   } 
	   } 
	 & & \\ 
 

    \theutterance \stepcounter{utterance}  

    & & & \multicolumn{2}{p{0.3\linewidth}}{\cellcolor[rgb]{0.95,0.95,0.95}{%
	\makecell[{{p{\linewidth}}}]{% 
	  \tt {\tiny [GM$|$GM]}  
	 SPOSTA: P, (1, 7) 
	  } 
	   } 
	   } 
	 & & \\ 
 

    \theutterance \stepcounter{utterance}  

    & & & \multicolumn{2}{p{0.3\linewidth}}{\cellcolor[rgb]{0.95,0.95,0.95}{%
	\makecell[{{p{\linewidth}}}]{% 
	  \tt {\tiny [GM$|$GM]}  
	 Spostato 'P' in (1,7) con successo. La tua griglia aggiornata appare così:\\ \tt ```\\ \tt  1234567\\ \tt ╔═╤═╤═╤═╗\\ \tt ║◌│◌│◌│A║ 1\\ \tt ╟─┴─┼─┤◌║ 2\\ \tt ║◌◌M│◌└─╢ 3\\ \tt ╟───┘◌S◌║ 4\\ \tt ║◌◌◌◌◌┌─╢ 5\\ \tt ║◌◌◌◌T│O║ 6\\ \tt ║PD◌◌◌│◌║ 7\\ \tt ╚═════╧═╝\\ \tt ```\\ \tt  
	  } 
	   } 
	   } 
	 & & \\ 
 

    \theutterance \stepcounter{utterance}  

    & & \multicolumn{4}{p{0.6\linewidth}}{\cellcolor[rgb]{0.9,0.9,0.9}{%
	\makecell[{{p{\linewidth}}}]{% 
	  \tt {\tiny [GM$\rangle$B]}  
	 Spostato 'P' in (1,7) con successo. La tua griglia aggiornata appare così:\\ \tt ```\\ \tt  1234567\\ \tt ╔═╤═╤═╤═╗\\ \tt ║◌│◌│◌│◌║ 1\\ \tt ╟─┴─┼─┤◌║ 2\\ \tt ║◌M◌│◌└─╢ 3\\ \tt ╟───┘◌◌◌║ 4\\ \tt ║O◌◌◌◌┌─╢ 5\\ \tt ║◌D◌◌A│S║ 6\\ \tt ║P◌◌◌T│◌║ 7\\ \tt ╚═════╧═╝\\ \tt ```\\ \tt You are currently playing round 3 of maximum 28.\\ \tt You have collectively accumulated 2 of 16 penalties.\\ \tt L'altro giocatore ha spostato un oggetto sulla sua griglia.\\ \tt \\ \tt \\ \tt Qual è il tuo prossimo comando? 
	  } 
	   } 
	   } 
	 & \\ 
 

    \theutterance \stepcounter{utterance}  

    & & & \multicolumn{4}{p{0.6\linewidth}}{\cellcolor[rgb]{1,0.85,0.72}{%
	\makecell[{{p{\linewidth}}}]{% 
	  \tt {\tiny [GM$\langle$B]}  
	 SPOSTA: T, (2,7) 
	  } 
	   } 
	   } 
	  \\ 
 

    \theutterance \stepcounter{utterance}  

    & & & \multicolumn{2}{p{0.3\linewidth}}{\cellcolor[rgb]{0.95,0.95,0.95}{%
	\makecell[{{p{\linewidth}}}]{% 
	  \tt {\tiny [GM$|$GM]}  
	 SPOSTA: T, (2,7) 
	  } 
	   } 
	   } 
	 & & \\ 
 

    \theutterance \stepcounter{utterance}  

    & & & \multicolumn{2}{p{0.3\linewidth}}{\cellcolor[rgb]{0.95,0.95,0.95}{%
	\makecell[{{p{\linewidth}}}]{% 
	  \tt {\tiny [GM$|$GM]}  
	 Spostato 'T' in (2,7) con successo. La tua griglia aggiornata appare così:\\ \tt ```\\ \tt  1234567\\ \tt ╔═╤═╤═╤═╗\\ \tt ║◌│◌│◌│◌║ 1\\ \tt ╟─┴─┼─┤◌║ 2\\ \tt ║◌M◌│◌└─╢ 3\\ \tt ╟───┘◌◌◌║ 4\\ \tt ║O◌◌◌◌┌─╢ 5\\ \tt ║◌D◌◌A│S║ 6\\ \tt ║PT◌◌◌│◌║ 7\\ \tt ╚═════╧═╝\\ \tt ```\\ \tt  
	  } 
	   } 
	   } 
	 & & \\ 
 

    \theutterance \stepcounter{utterance}  

    & & \multicolumn{4}{p{0.6\linewidth}}{\cellcolor[rgb]{0.9,0.9,0.9}{%
	\makecell[{{p{\linewidth}}}]{% 
	  \tt {\tiny [A$\langle$GM]}  
	 Spostato 'P' in (1,7) con successo. La tua griglia aggiornata appare così:\\ \tt ```\\ \tt  1234567\\ \tt ╔═╤═╤═╤═╗\\ \tt ║◌│◌│◌│A║ 1\\ \tt ╟─┴─┼─┤◌║ 2\\ \tt ║◌◌M│◌└─╢ 3\\ \tt ╟───┘◌S◌║ 4\\ \tt ║◌◌◌◌◌┌─╢ 5\\ \tt ║◌◌◌◌T│O║ 6\\ \tt ║PD◌◌◌│◌║ 7\\ \tt ╚═════╧═╝\\ \tt ```\\ \tt You are currently playing round 3 of maximum 28.\\ \tt You have collectively accumulated 2 of 16 penalties.\\ \tt L'altro giocatore ha spostato un oggetto sulla sua griglia.\\ \tt \\ \tt \\ \tt Qual è il tuo prossimo comando? 
	  } 
	   } 
	   } 
	 & \\ 
 

    \theutterance \stepcounter{utterance}  

    & \multicolumn{4}{p{0.6\linewidth}}{\cellcolor[rgb]{0.8,1,0.9}{%
	\makecell[{{p{\linewidth}}}]{% 
	  \tt {\tiny [A$\rangle$GM]}  
	 SPOSTA: T, (2, 7) 
	  } 
	   } 
	   } 
	 & & \\ 
 

    \theutterance \stepcounter{utterance}  

    & & & \multicolumn{2}{p{0.3\linewidth}}{\cellcolor[rgb]{0.95,0.95,0.95}{%
	\makecell[{{p{\linewidth}}}]{% 
	  \tt {\tiny [GM$|$GM]}  
	 SPOSTA: T, (2, 7) 
	  } 
	   } 
	   } 
	 & & \\ 
 

    \theutterance \stepcounter{utterance}  

    & & & \multicolumn{2}{p{0.3\linewidth}}{\cellcolor[rgb]{0.95,0.95,0.95}{%
	\makecell[{{p{\linewidth}}}]{% 
	  \tt {\tiny [GM$|$GM]}  
	 Penalità: (2,7) non è vuoto, ma contiene 'D'.\\ \tt \\ \tt You have collectively accumulated 3 of 16 penalties.\\ \tt Riprova! 
	  } 
	   } 
	   } 
	 & & \\ 
 

    \theutterance \stepcounter{utterance}  

    & & \multicolumn{4}{p{0.6\linewidth}}{\cellcolor[rgb]{0.9,0.9,0.9}{%
	\makecell[{{p{\linewidth}}}]{% 
	  \tt {\tiny [A$\langle$GM]}  
	 Penalità: (2,7) non è vuoto, ma contiene 'D'.\\ \tt \\ \tt You have collectively accumulated 3 of 16 penalties.\\ \tt Riprova! 
	  } 
	   } 
	   } 
	 & \\ 
 

    \theutterance \stepcounter{utterance}  

    & \multicolumn{4}{p{0.6\linewidth}}{\cellcolor[rgb]{0.8,1,0.9}{%
	\makecell[{{p{\linewidth}}}]{% 
	  \tt {\tiny [A$\rangle$GM]}  
	 COMUNICA: Ho notato che (2,7) ha D sulla mia griglia. Sposterò prima D per liberare spazio. Sposto D da (2,7) a (4,7) come pianificato. 
	  } 
	   } 
	   } 
	 & & \\ 
 

    \theutterance \stepcounter{utterance}  

    & & & \multicolumn{2}{p{0.3\linewidth}}{\cellcolor[rgb]{0.95,0.95,0.95}{%
	\makecell[{{p{\linewidth}}}]{% 
	  \tt {\tiny [GM$|$GM]}  
	 COMUNICA: Ho notato che (2,7) ha D sulla mia griglia. Sposterò prima D per liberare spazio. Sposto D da (2,7) a (4,7) come pianificato. 
	  } 
	   } 
	   } 
	 & & \\ 
 

    \theutterance \stepcounter{utterance}  

    & & \multicolumn{4}{p{0.6\linewidth}}{\cellcolor[rgb]{0.9,0.9,0.9}{%
	\makecell[{{p{\linewidth}}}]{% 
	  \tt {\tiny [GM$\rangle$B]}  
	 Spostato 'T' in (2,7) con successo. La tua griglia aggiornata appare così:\\ \tt ```\\ \tt  1234567\\ \tt ╔═╤═╤═╤═╗\\ \tt ║◌│◌│◌│◌║ 1\\ \tt ╟─┴─┼─┤◌║ 2\\ \tt ║◌M◌│◌└─╢ 3\\ \tt ╟───┘◌◌◌║ 4\\ \tt ║O◌◌◌◌┌─╢ 5\\ \tt ║◌D◌◌A│S║ 6\\ \tt ║PT◌◌◌│◌║ 7\\ \tt ╚═════╧═╝\\ \tt ```\\ \tt You are currently playing round 4 of maximum 28.\\ \tt You have collectively accumulated 3 of 16 penalties.\\ \tt L'altro giocatore ha inviato questo messaggio:\\ \tt \\ \tt \\ \tt "Ho notato che (2,7) ha D sulla mia griglia. Sposterò prima D per liberare spazio. Sposto D da (2,7) a (4,7) come pianificato."\\ \tt \\ \tt \\ \tt Qual è il tuo prossimo comando? 
	  } 
	   } 
	   } 
	 & \\ 
 

    \theutterance \stepcounter{utterance}  

    & & & \multicolumn{4}{p{0.6\linewidth}}{\cellcolor[rgb]{1,0.85,0.72}{%
	\makecell[{{p{\linewidth}}}]{% 
	  \tt {\tiny [GM$\langle$B]}  
	 SPOSTA: D, (4,7) 
	  } 
	   } 
	   } 
	  \\ 
 

    \theutterance \stepcounter{utterance}  

    & & & \multicolumn{2}{p{0.3\linewidth}}{\cellcolor[rgb]{0.95,0.95,0.95}{%
	\makecell[{{p{\linewidth}}}]{% 
	  \tt {\tiny [GM$|$GM]}  
	 SPOSTA: D, (4,7) 
	  } 
	   } 
	   } 
	 & & \\ 
 

    \theutterance \stepcounter{utterance}  

    & & & \multicolumn{2}{p{0.3\linewidth}}{\cellcolor[rgb]{0.95,0.95,0.95}{%
	\makecell[{{p{\linewidth}}}]{% 
	  \tt {\tiny [GM$|$GM]}  
	 Spostato 'D' in (4,7) con successo. La tua griglia aggiornata appare così:\\ \tt ```\\ \tt  1234567\\ \tt ╔═╤═╤═╤═╗\\ \tt ║◌│◌│◌│◌║ 1\\ \tt ╟─┴─┼─┤◌║ 2\\ \tt ║◌M◌│◌└─╢ 3\\ \tt ╟───┘◌◌◌║ 4\\ \tt ║O◌◌◌◌┌─╢ 5\\ \tt ║◌◌◌◌A│S║ 6\\ \tt ║PT◌D◌│◌║ 7\\ \tt ╚═════╧═╝\\ \tt ```\\ \tt  
	  } 
	   } 
	   } 
	 & & \\ 
 

    \theutterance \stepcounter{utterance}  

    & & \multicolumn{4}{p{0.6\linewidth}}{\cellcolor[rgb]{0.9,0.9,0.9}{%
	\makecell[{{p{\linewidth}}}]{% 
	  \tt {\tiny [A$\langle$GM]}  
	 Il tuo messaggio è stato inoltrato all'altro giocatore.\\ \tt \\ \tt \\ \tt You are currently playing round 4 of maximum 28.\\ \tt You have collectively accumulated 3 of 16 penalties.\\ \tt L'altro giocatore ha spostato un oggetto sulla sua griglia.\\ \tt \\ \tt \\ \tt Qual è il tuo prossimo comando? 
	  } 
	   } 
	   } 
	 & \\ 
 

    \theutterance \stepcounter{utterance}  

    & \multicolumn{4}{p{0.6\linewidth}}{\cellcolor[rgb]{0.8,1,0.9}{%
	\makecell[{{p{\linewidth}}}]{% 
	  \tt {\tiny [A$\rangle$GM]}  
	 SPOSTA: D, (4, 7) 
	  } 
	   } 
	   } 
	 & & \\ 
 

    \theutterance \stepcounter{utterance}  

    & & & \multicolumn{2}{p{0.3\linewidth}}{\cellcolor[rgb]{0.95,0.95,0.95}{%
	\makecell[{{p{\linewidth}}}]{% 
	  \tt {\tiny [GM$|$GM]}  
	 SPOSTA: D, (4, 7) 
	  } 
	   } 
	   } 
	 & & \\ 
 

    \theutterance \stepcounter{utterance}  

    & & & \multicolumn{2}{p{0.3\linewidth}}{\cellcolor[rgb]{0.95,0.95,0.95}{%
	\makecell[{{p{\linewidth}}}]{% 
	  \tt {\tiny [GM$|$GM]}  
	 Spostato 'D' in (4,7) con successo. La tua griglia aggiornata appare così:\\ \tt ```\\ \tt  1234567\\ \tt ╔═╤═╤═╤═╗\\ \tt ║◌│◌│◌│A║ 1\\ \tt ╟─┴─┼─┤◌║ 2\\ \tt ║◌◌M│◌└─╢ 3\\ \tt ╟───┘◌S◌║ 4\\ \tt ║◌◌◌◌◌┌─╢ 5\\ \tt ║◌◌◌◌T│O║ 6\\ \tt ║P◌◌D◌│◌║ 7\\ \tt ╚═════╧═╝\\ \tt ```\\ \tt  
	  } 
	   } 
	   } 
	 & & \\ 
 

    \theutterance \stepcounter{utterance}  

    & & \multicolumn{4}{p{0.6\linewidth}}{\cellcolor[rgb]{0.9,0.9,0.9}{%
	\makecell[{{p{\linewidth}}}]{% 
	  \tt {\tiny [GM$\rangle$B]}  
	 Spostato 'D' in (4,7) con successo. La tua griglia aggiornata appare così:\\ \tt ```\\ \tt  1234567\\ \tt ╔═╤═╤═╤═╗\\ \tt ║◌│◌│◌│◌║ 1\\ \tt ╟─┴─┼─┤◌║ 2\\ \tt ║◌M◌│◌└─╢ 3\\ \tt ╟───┘◌◌◌║ 4\\ \tt ║O◌◌◌◌┌─╢ 5\\ \tt ║◌◌◌◌A│S║ 6\\ \tt ║PT◌D◌│◌║ 7\\ \tt ╚═════╧═╝\\ \tt ```\\ \tt You are currently playing round 5 of maximum 28.\\ \tt You have collectively accumulated 3 of 16 penalties.\\ \tt L'altro giocatore ha spostato un oggetto sulla sua griglia.\\ \tt \\ \tt \\ \tt Qual è il tuo prossimo comando? 
	  } 
	   } 
	   } 
	 & \\ 
 

    \theutterance \stepcounter{utterance}  

    & & & \multicolumn{4}{p{0.6\linewidth}}{\cellcolor[rgb]{1,0.85,0.72}{%
	\makecell[{{p{\linewidth}}}]{% 
	  \tt {\tiny [GM$\langle$B]}  
	 SPOSTA: A, (5,7) 
	  } 
	   } 
	   } 
	  \\ 
 

    \theutterance \stepcounter{utterance}  

    & & & \multicolumn{2}{p{0.3\linewidth}}{\cellcolor[rgb]{0.95,0.95,0.95}{%
	\makecell[{{p{\linewidth}}}]{% 
	  \tt {\tiny [GM$|$GM]}  
	 SPOSTA: A, (5,7) 
	  } 
	   } 
	   } 
	 & & \\ 
 

    \theutterance \stepcounter{utterance}  

    & & & \multicolumn{2}{p{0.3\linewidth}}{\cellcolor[rgb]{0.95,0.95,0.95}{%
	\makecell[{{p{\linewidth}}}]{% 
	  \tt {\tiny [GM$|$GM]}  
	 Spostato 'A' in (5,7) con successo. La tua griglia aggiornata appare così:\\ \tt ```\\ \tt  1234567\\ \tt ╔═╤═╤═╤═╗\\ \tt ║◌│◌│◌│◌║ 1\\ \tt ╟─┴─┼─┤◌║ 2\\ \tt ║◌M◌│◌└─╢ 3\\ \tt ╟───┘◌◌◌║ 4\\ \tt ║O◌◌◌◌┌─╢ 5\\ \tt ║◌◌◌◌◌│S║ 6\\ \tt ║PT◌DA│◌║ 7\\ \tt ╚═════╧═╝\\ \tt ```\\ \tt  
	  } 
	   } 
	   } 
	 & & \\ 
 

    \theutterance \stepcounter{utterance}  

    & & \multicolumn{4}{p{0.6\linewidth}}{\cellcolor[rgb]{0.9,0.9,0.9}{%
	\makecell[{{p{\linewidth}}}]{% 
	  \tt {\tiny [A$\langle$GM]}  
	 Spostato 'D' in (4,7) con successo. La tua griglia aggiornata appare così:\\ \tt ```\\ \tt  1234567\\ \tt ╔═╤═╤═╤═╗\\ \tt ║◌│◌│◌│A║ 1\\ \tt ╟─┴─┼─┤◌║ 2\\ \tt ║◌◌M│◌└─╢ 3\\ \tt ╟───┘◌S◌║ 4\\ \tt ║◌◌◌◌◌┌─╢ 5\\ \tt ║◌◌◌◌T│O║ 6\\ \tt ║P◌◌D◌│◌║ 7\\ \tt ╚═════╧═╝\\ \tt ```\\ \tt You are currently playing round 5 of maximum 28.\\ \tt You have collectively accumulated 3 of 16 penalties.\\ \tt L'altro giocatore ha spostato un oggetto sulla sua griglia.\\ \tt \\ \tt \\ \tt Qual è il tuo prossimo comando? 
	  } 
	   } 
	   } 
	 & \\ 
 

    \theutterance \stepcounter{utterance}  

    & \multicolumn{4}{p{0.6\linewidth}}{\cellcolor[rgb]{0.8,1,0.9}{%
	\makecell[{{p{\linewidth}}}]{% 
	  \tt {\tiny [A$\rangle$GM]}  
	 SPOSTA: T, (2, 7) 
	  } 
	   } 
	   } 
	 & & \\ 
 

    \theutterance \stepcounter{utterance}  

    & & & \multicolumn{2}{p{0.3\linewidth}}{\cellcolor[rgb]{0.95,0.95,0.95}{%
	\makecell[{{p{\linewidth}}}]{% 
	  \tt {\tiny [GM$|$GM]}  
	 SPOSTA: T, (2, 7) 
	  } 
	   } 
	   } 
	 & & \\ 
 

    \theutterance \stepcounter{utterance}  

    & & & \multicolumn{2}{p{0.3\linewidth}}{\cellcolor[rgb]{0.95,0.95,0.95}{%
	\makecell[{{p{\linewidth}}}]{% 
	  \tt {\tiny [GM$|$GM]}  
	 Spostato 'T' in (2,7) con successo. La tua griglia aggiornata appare così:\\ \tt ```\\ \tt  1234567\\ \tt ╔═╤═╤═╤═╗\\ \tt ║◌│◌│◌│A║ 1\\ \tt ╟─┴─┼─┤◌║ 2\\ \tt ║◌◌M│◌└─╢ 3\\ \tt ╟───┘◌S◌║ 4\\ \tt ║◌◌◌◌◌┌─╢ 5\\ \tt ║◌◌◌◌◌│O║ 6\\ \tt ║PT◌D◌│◌║ 7\\ \tt ╚═════╧═╝\\ \tt ```\\ \tt  
	  } 
	   } 
	   } 
	 & & \\ 
 

    \theutterance \stepcounter{utterance}  

    & & \multicolumn{4}{p{0.6\linewidth}}{\cellcolor[rgb]{0.9,0.9,0.9}{%
	\makecell[{{p{\linewidth}}}]{% 
	  \tt {\tiny [GM$\rangle$B]}  
	 Spostato 'A' in (5,7) con successo. La tua griglia aggiornata appare così:\\ \tt ```\\ \tt  1234567\\ \tt ╔═╤═╤═╤═╗\\ \tt ║◌│◌│◌│◌║ 1\\ \tt ╟─┴─┼─┤◌║ 2\\ \tt ║◌M◌│◌└─╢ 3\\ \tt ╟───┘◌◌◌║ 4\\ \tt ║O◌◌◌◌┌─╢ 5\\ \tt ║◌◌◌◌◌│S║ 6\\ \tt ║PT◌DA│◌║ 7\\ \tt ╚═════╧═╝\\ \tt ```\\ \tt You are currently playing round 6 of maximum 28.\\ \tt You have collectively accumulated 3 of 16 penalties.\\ \tt L'altro giocatore ha spostato un oggetto sulla sua griglia.\\ \tt \\ \tt \\ \tt Qual è il tuo prossimo comando? 
	  } 
	   } 
	   } 
	 & \\ 
 

    \theutterance \stepcounter{utterance}  

    & & & \multicolumn{4}{p{0.6\linewidth}}{\cellcolor[rgb]{1,0.85,0.72}{%
	\makecell[{{p{\linewidth}}}]{% 
	  \tt {\tiny [GM$\langle$B]}  
	 SPOSTA: S, (3,7) 
	  } 
	   } 
	   } 
	  \\ 
 

    \theutterance \stepcounter{utterance}  

    & & & \multicolumn{2}{p{0.3\linewidth}}{\cellcolor[rgb]{0.95,0.95,0.95}{%
	\makecell[{{p{\linewidth}}}]{% 
	  \tt {\tiny [GM$|$GM]}  
	 SPOSTA: S, (3,7) 
	  } 
	   } 
	   } 
	 & & \\ 
 

    \theutterance \stepcounter{utterance}  

    & & & \multicolumn{2}{p{0.3\linewidth}}{\cellcolor[rgb]{0.95,0.95,0.95}{%
	\makecell[{{p{\linewidth}}}]{% 
	  \tt {\tiny [GM$|$GM]}  
	 Spostato 'S' in (3,7) con successo. La tua griglia aggiornata appare così:\\ \tt ```\\ \tt  1234567\\ \tt ╔═╤═╤═╤═╗\\ \tt ║◌│◌│◌│◌║ 1\\ \tt ╟─┴─┼─┤◌║ 2\\ \tt ║◌M◌│◌└─╢ 3\\ \tt ╟───┘◌◌◌║ 4\\ \tt ║O◌◌◌◌┌─╢ 5\\ \tt ║◌◌◌◌◌│◌║ 6\\ \tt ║PTSDA│◌║ 7\\ \tt ╚═════╧═╝\\ \tt ```\\ \tt  
	  } 
	   } 
	   } 
	 & & \\ 
 

    \theutterance \stepcounter{utterance}  

    & & \multicolumn{4}{p{0.6\linewidth}}{\cellcolor[rgb]{0.9,0.9,0.9}{%
	\makecell[{{p{\linewidth}}}]{% 
	  \tt {\tiny [A$\langle$GM]}  
	 Spostato 'T' in (2,7) con successo. La tua griglia aggiornata appare così:\\ \tt ```\\ \tt  1234567\\ \tt ╔═╤═╤═╤═╗\\ \tt ║◌│◌│◌│A║ 1\\ \tt ╟─┴─┼─┤◌║ 2\\ \tt ║◌◌M│◌└─╢ 3\\ \tt ╟───┘◌S◌║ 4\\ \tt ║◌◌◌◌◌┌─╢ 5\\ \tt ║◌◌◌◌◌│O║ 6\\ \tt ║PT◌D◌│◌║ 7\\ \tt ╚═════╧═╝\\ \tt ```\\ \tt You are currently playing round 6 of maximum 28.\\ \tt You have collectively accumulated 3 of 16 penalties.\\ \tt L'altro giocatore ha spostato un oggetto sulla sua griglia.\\ \tt \\ \tt \\ \tt Qual è il tuo prossimo comando? 
	  } 
	   } 
	   } 
	 & \\ 
 

    \theutterance \stepcounter{utterance}  

    & \multicolumn{4}{p{0.6\linewidth}}{\cellcolor[rgb]{0.8,1,0.9}{%
	\makecell[{{p{\linewidth}}}]{% 
	  \tt {\tiny [A$\rangle$GM]}  
	 COMUNICA: Ho spostato P(1,7), T(2,7), D(4,7). Ora sposterò S da (5,4) a (3,7). Conferma i tuoi progressi. 
	  } 
	   } 
	   } 
	 & & \\ 
 

    \theutterance \stepcounter{utterance}  

    & & & \multicolumn{2}{p{0.3\linewidth}}{\cellcolor[rgb]{0.95,0.95,0.95}{%
	\makecell[{{p{\linewidth}}}]{% 
	  \tt {\tiny [GM$|$GM]}  
	 COMUNICA: Ho spostato P(1,7), T(2,7), D(4,7). Ora sposterò S da (5,4) a (3,7). Conferma i tuoi progressi. 
	  } 
	   } 
	   } 
	 & & \\ 
 

    \theutterance \stepcounter{utterance}  

    & & \multicolumn{4}{p{0.6\linewidth}}{\cellcolor[rgb]{0.9,0.9,0.9}{%
	\makecell[{{p{\linewidth}}}]{% 
	  \tt {\tiny [GM$\rangle$B]}  
	 Spostato 'S' in (3,7) con successo. La tua griglia aggiornata appare così:\\ \tt ```\\ \tt  1234567\\ \tt ╔═╤═╤═╤═╗\\ \tt ║◌│◌│◌│◌║ 1\\ \tt ╟─┴─┼─┤◌║ 2\\ \tt ║◌M◌│◌└─╢ 3\\ \tt ╟───┘◌◌◌║ 4\\ \tt ║O◌◌◌◌┌─╢ 5\\ \tt ║◌◌◌◌◌│◌║ 6\\ \tt ║PTSDA│◌║ 7\\ \tt ╚═════╧═╝\\ \tt ```\\ \tt You are currently playing round 7 of maximum 28.\\ \tt You have collectively accumulated 3 of 16 penalties.\\ \tt L'altro giocatore ha inviato questo messaggio:\\ \tt \\ \tt \\ \tt "Ho spostato P(1,7), T(2,7), D(4,7). Ora sposterò S da (5,4) a (3,7). Conferma i tuoi progressi."\\ \tt \\ \tt \\ \tt Qual è il tuo prossimo comando? 
	  } 
	   } 
	   } 
	 & \\ 
 

    \theutterance \stepcounter{utterance}  

    & & & \multicolumn{4}{p{0.6\linewidth}}{\cellcolor[rgb]{1,0.85,0.72}{%
	\makecell[{{p{\linewidth}}}]{% 
	  \tt {\tiny [GM$\langle$B]}  
	 SPOSTA: M, (6,7) 
	  } 
	   } 
	   } 
	  \\ 
 

    \theutterance \stepcounter{utterance}  

    & & & \multicolumn{2}{p{0.3\linewidth}}{\cellcolor[rgb]{0.95,0.95,0.95}{%
	\makecell[{{p{\linewidth}}}]{% 
	  \tt {\tiny [GM$|$GM]}  
	 SPOSTA: M, (6,7) 
	  } 
	   } 
	   } 
	 & & \\ 
 

    \theutterance \stepcounter{utterance}  

    & & & \multicolumn{2}{p{0.3\linewidth}}{\cellcolor[rgb]{0.95,0.95,0.95}{%
	\makecell[{{p{\linewidth}}}]{% 
	  \tt {\tiny [GM$|$GM]}  
	 Penalità: (6,7) non è vuoto, ma contiene '│'.\\ \tt \\ \tt You have collectively accumulated 4 of 16 penalties.\\ \tt Riprova! 
	  } 
	   } 
	   } 
	 & & \\ 
 

    \theutterance \stepcounter{utterance}  

    & & \multicolumn{4}{p{0.6\linewidth}}{\cellcolor[rgb]{0.9,0.9,0.9}{%
	\makecell[{{p{\linewidth}}}]{% 
	  \tt {\tiny [GM$\rangle$B]}  
	 Penalità: (6,7) non è vuoto, ma contiene '│'.\\ \tt \\ \tt You have collectively accumulated 4 of 16 penalties.\\ \tt Riprova! 
	  } 
	   } 
	   } 
	 & \\ 
 

    \theutterance \stepcounter{utterance}  

    & & & \multicolumn{4}{p{0.6\linewidth}}{\cellcolor[rgb]{1,0.85,0.72}{%
	\makecell[{{p{\linewidth}}}]{% 
	  \tt {\tiny [GM$\langle$B]}  
	 La cella (6,7) contiene il simbolo '│' che non è considerato vuoto ('◌'). Devo trovare una cella libera nella riga 7. Osservando la griglia, vedo che (7,7) è libera (contiene '◌').\\ \tt \\ \tt Prima sposto O in (7,7) per liberare lo spazio per M, poi sposterò M in (6,7) dopo aver controllato che sia libero.\\ \tt \\ \tt SPOSTA: O, (7,7) 
	  } 
	   } 
	   } 
	  \\ 
 

    \theutterance \stepcounter{utterance}  

    & & & \multicolumn{2}{p{0.3\linewidth}}{\cellcolor[rgb]{0.95,0.95,0.95}{%
	\makecell[{{p{\linewidth}}}]{% 
	  \tt {\tiny [GM$|$GM]}  
	 La cella (6,7) contiene il simbolo '│' che non è considerato vuoto ('◌'). Devo trovare una cella libera nella riga 7. Osservando la griglia, vedo che (7,7) è libera (contiene '◌').\\ \tt \\ \tt Prima sposto O in (7,7) per liberare lo spazio per M, poi sposterò M in (6,7) dopo aver controllato che sia libero.\\ \tt \\ \tt SPOSTA: O, (7,7) 
	  } 
	   } 
	   } 
	 & & \\ 
 

    \theutterance \stepcounter{utterance}  

    & & & \multicolumn{2}{p{0.3\linewidth}}{\cellcolor[rgb]{0.95,0.95,0.95}{%
	\makecell[{{p{\linewidth}}}]{% 
	  \tt {\tiny [GM$|$GM]}  
	 Invalid format: head is not empty: 'La cella (6,7) contiene il simbolo '│' che non è considerato vuoto ('◌'). Devo trovare una cella libera nella riga 7. Osservando la griglia, vedo che (7,7) è libera (contiene '◌').\\ \tt \\ \tt Prima sposto O in (7,7) per liberare lo spazio per M, poi sposterò M in (6,7) dopo aver controllato che sia libero.\\ \tt \\ \tt ' 
	  } 
	   } 
	   } 
	 & & \\ 
 

    \theutterance \stepcounter{utterance}  

    & & \multicolumn{4}{p{0.6\linewidth}}{\cellcolor[rgb]{0.9,0.9,0.9}{%
	\makecell[{{p{\linewidth}}}]{% 
	  \tt {\tiny [GM$\rangle$B]}  
	 Penalità: Il tuo messaggio non deve contenere nulla prima del comando!\\ \tt \\ \tt Assicurati che la tua risposta contenga solo `COMUNICA: <MESSAGE>` o `SPOSTA: <OBJECT>, (<X>, <Y>)`, e nient'altro!\\ \tt \\ \tt You have collectively accumulated 5 of 16 penalties.\\ \tt Riprova! 
	  } 
	   } 
	   } 
	 & \\ 
 

    \theutterance \stepcounter{utterance}  

    & & & \multicolumn{4}{p{0.6\linewidth}}{\cellcolor[rgb]{1,0.85,0.72}{%
	\makecell[{{p{\linewidth}}}]{% 
	  \tt {\tiny [GM$\langle$B]}  
	 SPOSTA: O, (7,7) 
	  } 
	   } 
	   } 
	  \\ 
 

    \theutterance \stepcounter{utterance}  

    & & & \multicolumn{2}{p{0.3\linewidth}}{\cellcolor[rgb]{0.95,0.95,0.95}{%
	\makecell[{{p{\linewidth}}}]{% 
	  \tt {\tiny [GM$|$GM]}  
	 SPOSTA: O, (7,7) 
	  } 
	   } 
	   } 
	 & & \\ 
 

    \theutterance \stepcounter{utterance}  

    & & & \multicolumn{2}{p{0.3\linewidth}}{\cellcolor[rgb]{0.95,0.95,0.95}{%
	\makecell[{{p{\linewidth}}}]{% 
	  \tt {\tiny [GM$|$GM]}  
	 Spostato 'O' in (7,7) con successo. La tua griglia aggiornata appare così:\\ \tt ```\\ \tt  1234567\\ \tt ╔═╤═╤═╤═╗\\ \tt ║◌│◌│◌│◌║ 1\\ \tt ╟─┴─┼─┤◌║ 2\\ \tt ║◌M◌│◌└─╢ 3\\ \tt ╟───┘◌◌◌║ 4\\ \tt ║◌◌◌◌◌┌─╢ 5\\ \tt ║◌◌◌◌◌│◌║ 6\\ \tt ║PTSDA│O║ 7\\ \tt ╚═════╧═╝\\ \tt ```\\ \tt  
	  } 
	   } 
	   } 
	 & & \\ 
 

    \theutterance \stepcounter{utterance}  

    & & \multicolumn{4}{p{0.6\linewidth}}{\cellcolor[rgb]{0.9,0.9,0.9}{%
	\makecell[{{p{\linewidth}}}]{% 
	  \tt {\tiny [A$\langle$GM]}  
	 Il tuo messaggio è stato inoltrato all'altro giocatore.\\ \tt \\ \tt \\ \tt You are currently playing round 7 of maximum 28.\\ \tt You have collectively accumulated 5 of 16 penalties.\\ \tt L'altro giocatore ha spostato un oggetto sulla sua griglia.\\ \tt \\ \tt \\ \tt Qual è il tuo prossimo comando? 
	  } 
	   } 
	   } 
	 & \\ 
 

    \theutterance \stepcounter{utterance}  

    & \multicolumn{4}{p{0.6\linewidth}}{\cellcolor[rgb]{0.8,1,0.9}{%
	\makecell[{{p{\linewidth}}}]{% 
	  \tt {\tiny [A$\rangle$GM]}  
	 SPOSTA: S, (3, 7) 
	  } 
	   } 
	   } 
	 & & \\ 
 

    \theutterance \stepcounter{utterance}  

    & & & \multicolumn{2}{p{0.3\linewidth}}{\cellcolor[rgb]{0.95,0.95,0.95}{%
	\makecell[{{p{\linewidth}}}]{% 
	  \tt {\tiny [GM$|$GM]}  
	 SPOSTA: S, (3, 7) 
	  } 
	   } 
	   } 
	 & & \\ 
 

    \theutterance \stepcounter{utterance}  

    & & & \multicolumn{2}{p{0.3\linewidth}}{\cellcolor[rgb]{0.95,0.95,0.95}{%
	\makecell[{{p{\linewidth}}}]{% 
	  \tt {\tiny [GM$|$GM]}  
	 Spostato 'S' in (3,7) con successo. La tua griglia aggiornata appare così:\\ \tt ```\\ \tt  1234567\\ \tt ╔═╤═╤═╤═╗\\ \tt ║◌│◌│◌│A║ 1\\ \tt ╟─┴─┼─┤◌║ 2\\ \tt ║◌◌M│◌└─╢ 3\\ \tt ╟───┘◌◌◌║ 4\\ \tt ║◌◌◌◌◌┌─╢ 5\\ \tt ║◌◌◌◌◌│O║ 6\\ \tt ║PTSD◌│◌║ 7\\ \tt ╚═════╧═╝\\ \tt ```\\ \tt  
	  } 
	   } 
	   } 
	 & & \\ 
 

    \theutterance \stepcounter{utterance}  

    & & \multicolumn{4}{p{0.6\linewidth}}{\cellcolor[rgb]{0.9,0.9,0.9}{%
	\makecell[{{p{\linewidth}}}]{% 
	  \tt {\tiny [GM$\rangle$B]}  
	 Spostato 'O' in (7,7) con successo. La tua griglia aggiornata appare così:\\ \tt ```\\ \tt  1234567\\ \tt ╔═╤═╤═╤═╗\\ \tt ║◌│◌│◌│◌║ 1\\ \tt ╟─┴─┼─┤◌║ 2\\ \tt ║◌M◌│◌└─╢ 3\\ \tt ╟───┘◌◌◌║ 4\\ \tt ║◌◌◌◌◌┌─╢ 5\\ \tt ║◌◌◌◌◌│◌║ 6\\ \tt ║PTSDA│O║ 7\\ \tt ╚═════╧═╝\\ \tt ```\\ \tt You are currently playing round 8 of maximum 28.\\ \tt You have collectively accumulated 5 of 16 penalties.\\ \tt L'altro giocatore ha spostato un oggetto sulla sua griglia.\\ \tt \\ \tt \\ \tt Qual è il tuo prossimo comando? 
	  } 
	   } 
	   } 
	 & \\ 
 

    \theutterance \stepcounter{utterance}  

    & & & \multicolumn{4}{p{0.6\linewidth}}{\cellcolor[rgb]{1,0.85,0.72}{%
	\makecell[{{p{\linewidth}}}]{% 
	  \tt {\tiny [GM$\langle$B]}  
	 SPOSTA: M, (6,7) 
	  } 
	   } 
	   } 
	  \\ 
 

    \theutterance \stepcounter{utterance}  

    & & & \multicolumn{2}{p{0.3\linewidth}}{\cellcolor[rgb]{0.95,0.95,0.95}{%
	\makecell[{{p{\linewidth}}}]{% 
	  \tt {\tiny [GM$|$GM]}  
	 SPOSTA: M, (6,7) 
	  } 
	   } 
	   } 
	 & & \\ 
 

    \theutterance \stepcounter{utterance}  

    & & & \multicolumn{2}{p{0.3\linewidth}}{\cellcolor[rgb]{0.95,0.95,0.95}{%
	\makecell[{{p{\linewidth}}}]{% 
	  \tt {\tiny [GM$|$GM]}  
	 Penalità: (6,7) non è vuoto, ma contiene '│'.\\ \tt \\ \tt You have collectively accumulated 6 of 16 penalties.\\ \tt Riprova! 
	  } 
	   } 
	   } 
	 & & \\ 
 

    \theutterance \stepcounter{utterance}  

    & & \multicolumn{4}{p{0.6\linewidth}}{\cellcolor[rgb]{0.9,0.9,0.9}{%
	\makecell[{{p{\linewidth}}}]{% 
	  \tt {\tiny [GM$\rangle$B]}  
	 Penalità: (6,7) non è vuoto, ma contiene '│'.\\ \tt \\ \tt You have collectively accumulated 6 of 16 penalties.\\ \tt Riprova! 
	  } 
	   } 
	   } 
	 & \\ 
 

    \theutterance \stepcounter{utterance}  

    & & & \multicolumn{4}{p{0.6\linewidth}}{\cellcolor[rgb]{1,0.85,0.72}{%
	\makecell[{{p{\linewidth}}}]{% 
	  \tt {\tiny [GM$\langle$B]}  
	 Osservando la griglia, vedo che la cella (6,7) contiene il carattere '│', che non è una cella vuota. La riga 7 ha solo le colonne 1-5 e 7 come celle valide (vuote) per gli oggetti, poiché le colonne 6 e parte della 7 sono occupate da bordi.\\ \tt \\ \tt Dalla griglia:\\ \tt ```\\ \tt ║PTSDA│O║ 7\\ \tt ```\\ \tt Le posizioni sono:\\ \tt - Colonna 1: P\\ \tt - Colonna 2: T\\ \tt - Colonna 3: S\\ \tt - Colonna 4: D\\ \tt - Colonna 5: A\\ \tt - Colonna 6: │ (non vuoto)\\ \tt - Colonna 7: O\\ \tt \\ \tt Quindi, la posizione (6,7) non è valida. La proposta originale era M(6,7), ma ora vedo che (6,7) non è disponibile. Devo comunicare con l'altro giocatore per aggiustare la strategia.\\ \tt \\ \tt Invierò un messaggio per proporre una modifica.\\ \tt \\ \tt COMUNICA: (6,7) non è vuoto sulla mia griglia. Propongo di mettere M in (5,7) e spostare A in (6,7)? Ma (6,7) ha '│'. Alternative: usare riga 6 o altra riga? 
	  } 
	   } 
	   } 
	  \\ 
 

    \theutterance \stepcounter{utterance}  

    & & & \multicolumn{2}{p{0.3\linewidth}}{\cellcolor[rgb]{0.95,0.95,0.95}{%
	\makecell[{{p{\linewidth}}}]{% 
	  \tt {\tiny [GM$|$GM]}  
	 Osservando la griglia, vedo che la cella (6,7) contiene il carattere '│', che non è una cella vuota. La riga 7 ha solo le colonne 1-5 e 7 come celle valide (vuote) per gli oggetti, poiché le colonne 6 e parte della 7 sono occupate da bordi.\\ \tt \\ \tt Dalla griglia:\\ \tt ```\\ \tt ║PTSDA│O║ 7\\ \tt ```\\ \tt Le posizioni sono:\\ \tt - Colonna 1: P\\ \tt - Colonna 2: T\\ \tt - Colonna 3: S\\ \tt - Colonna 4: D\\ \tt - Colonna 5: A\\ \tt - Colonna 6: │ (non vuoto)\\ \tt - Colonna 7: O\\ \tt \\ \tt Quindi, la posizione (6,7) non è valida. La proposta originale era M(6,7), ma ora vedo che (6,7) non è disponibile. Devo comunicare con l'altro giocatore per aggiustare la strategia.\\ \tt \\ \tt Invierò un messaggio per proporre una modifica.\\ \tt \\ \tt COMUNICA: (6,7) non è vuoto sulla mia griglia. Propongo di mettere M in (5,7) e spostare A in (6,7)? Ma (6,7) ha '│'. Alternative: usare riga 6 o altra riga? 
	  } 
	   } 
	   } 
	 & & \\ 
 

    \theutterance \stepcounter{utterance}  

    & & & \multicolumn{2}{p{0.3\linewidth}}{\cellcolor[rgb]{0.95,0.95,0.95}{%
	\makecell[{{p{\linewidth}}}]{% 
	  \tt {\tiny [GM$|$GM]}  
	 Invalid format: head is not empty: 'Osservando la griglia, vedo che la cella (6,7) contiene il carattere '│', che non è una cella vuota. La riga 7 ha solo le colonne 1-5 e 7 come celle valide (vuote) per gli oggetti, poiché le colonne 6 e parte della 7 sono occupate da bordi.\\ \tt \\ \tt Dalla griglia:\\ \tt \\ \tt ║PTSDA│O║ 7\\ \tt \\ \tt Le posizioni sono:\\ \tt - Colonna 1: P\\ \tt - Colonna 2: T\\ \tt - Colonna 3: S\\ \tt - Colonna 4: D\\ \tt - Colonna 5: A\\ \tt - Colonna 6: │ (non vuoto)\\ \tt - Colonna 7: O\\ \tt \\ \tt Quindi, la posizione (6,7) non è valida. La proposta originale era M(6,7), ma ora vedo che (6,7) non è disponibile. Devo comunicare con l'altro giocatore per aggiustare la strategia.\\ \tt \\ \tt Invierò un messaggio per proporre una modifica.\\ \tt \\ \tt ' 
	  } 
	   } 
	   } 
	 & & \\ 
 

    \theutterance \stepcounter{utterance}  

    & & \multicolumn{4}{p{0.6\linewidth}}{\cellcolor[rgb]{0.9,0.9,0.9}{%
	\makecell[{{p{\linewidth}}}]{% 
	  \tt {\tiny [GM$\rangle$B]}  
	 Penalità: Il tuo messaggio non deve contenere nulla prima del comando!\\ \tt \\ \tt Assicurati che la tua risposta contenga solo `COMUNICA: <MESSAGE>` o `SPOSTA: <OBJECT>, (<X>, <Y>)`, e nient'altro!\\ \tt \\ \tt You have collectively accumulated 7 of 16 penalties.\\ \tt Riprova! 
	  } 
	   } 
	   } 
	 & \\ 
 

    \theutterance \stepcounter{utterance}  

    & & & \multicolumn{4}{p{0.6\linewidth}}{\cellcolor[rgb]{1,0.85,0.72}{%
	\makecell[{{p{\linewidth}}}]{% 
	  \tt {\tiny [GM$\langle$B]}  
	 COMUNICA: (6,7) non è vuoto sulla mia griglia. Propongo di mettere M in (5,7) e spostare A in (6,7)? Ma (6,7) ha '│'. Alternative: usare riga 6 o altra riga? 
	  } 
	   } 
	   } 
	  \\ 
 

    \theutterance \stepcounter{utterance}  

    & & & \multicolumn{2}{p{0.3\linewidth}}{\cellcolor[rgb]{0.95,0.95,0.95}{%
	\makecell[{{p{\linewidth}}}]{% 
	  \tt {\tiny [GM$|$GM]}  
	 COMUNICA: (6,7) non è vuoto sulla mia griglia. Propongo di mettere M in (5,7) e spostare A in (6,7)? Ma (6,7) ha '│'. Alternative: usare riga 6 o altra riga? 
	  } 
	   } 
	   } 
	 & & \\ 
 

    \theutterance \stepcounter{utterance}  

    & & \multicolumn{4}{p{0.6\linewidth}}{\cellcolor[rgb]{0.9,0.9,0.9}{%
	\makecell[{{p{\linewidth}}}]{% 
	  \tt {\tiny [A$\langle$GM]}  
	 Spostato 'S' in (3,7) con successo. La tua griglia aggiornata appare così:\\ \tt ```\\ \tt  1234567\\ \tt ╔═╤═╤═╤═╗\\ \tt ║◌│◌│◌│A║ 1\\ \tt ╟─┴─┼─┤◌║ 2\\ \tt ║◌◌M│◌└─╢ 3\\ \tt ╟───┘◌◌◌║ 4\\ \tt ║◌◌◌◌◌┌─╢ 5\\ \tt ║◌◌◌◌◌│O║ 6\\ \tt ║PTSD◌│◌║ 7\\ \tt ╚═════╧═╝\\ \tt ```\\ \tt You are currently playing round 8 of maximum 28.\\ \tt You have collectively accumulated 7 of 16 penalties.\\ \tt L'altro giocatore ha inviato questo messaggio:\\ \tt \\ \tt \\ \tt "(6,7) non è vuoto sulla mia griglia. Propongo di mettere M in (5,7) e spostare A in (6,7)? Ma (6,7) ha '│'. Alternative: usare riga 6 o altra riga?"\\ \tt \\ \tt \\ \tt Qual è il tuo prossimo comando? 
	  } 
	   } 
	   } 
	 & \\ 
 

    \theutterance \stepcounter{utterance}  

    & \multicolumn{4}{p{0.6\linewidth}}{\cellcolor[rgb]{0.8,1,0.9}{%
	\makecell[{{p{\linewidth}}}]{% 
	  \tt {\tiny [A$\rangle$GM]}  
	 COMUNICA: Vedo che (6,7) ha un bordo '│' sulla mia griglia, quindi non è una cella valida. Controllo la mia griglia: le celle valide in y=7 sono solo x=1,2,3,4,5,7 (x=6 è bordo). Quindi propongo di modificare la configurazione obiettivo: usiamo y=7 per x=1,2,3,4,5,7 e mettiamo M in (5,7) e A in (7,7)? O forse usiamo la riga 6? Le celle valide in y=6: x=1,2,3,4,5,7 (x=6 è bordo). Quindi possiamo usare:\\ \tt P(1,6), T(2,6), S(3,6), D(4,6), A(5,6), M(7,6)? Ma O è già in (7,6). Quindi forse spostare O altrove.\\ \tt Alternative: usare riga 5? Le celle valide in y=5: x=1,2,3,4,5 (x=6 è bordo). Non abbastanza spazio.\\ \tt Meglio usare riga 7 con x=1,2,3,4,5,7: quindi P(1,7), T(2,7), S(3,7), D(4,7), A(5,7), O(7,7) e M in (6,6) o altro? Ma (6,6) è bordo.\\ \tt Propongo: P(1,7), T(2,7), S(3,7), D(4,7), A(5,7), O(7,7) e M in (7,5) o (1,6)? Dobbiamo trovare una cella libera per M.\\ \tt Sulla mia griglia, (7,5) è libero? y=5: "║◌◌◌◌◌┌─╢" so x=1-5 sono ◌, x=6 è bordo, x=7? Non è visibile, ma probabilmente è bordo. In realtà, la griglia ha x=1-7, ma i bordi occupano alcune celle.\\ \tt Per evitare confusione, propongo di usare la riga 6 per alcuni oggetti: ad esempio, mettere M in (1,6) e O in (7,6) (già lì), e gli altri in riga 7.\\ \tt Quindi obiettivo modificato:\\ \tt P(1,7), T(2,7), S(3,7), D(4,7), A(5,7), M(1,6), O(7,6). Confermi? 
	  } 
	   } 
	   } 
	 & & \\ 
 

    \theutterance \stepcounter{utterance}  

    & & & \multicolumn{2}{p{0.3\linewidth}}{\cellcolor[rgb]{0.95,0.95,0.95}{%
	\makecell[{{p{\linewidth}}}]{% 
	  \tt {\tiny [GM$|$GM]}  
	 COMUNICA: Vedo che (6,7) ha un bordo '│' sulla mia griglia, quindi non è una cella valida. Controllo la mia griglia: le celle valide in y=7 sono solo x=1,2,3,4,5,7 (x=6 è bordo). Quindi propongo di modificare la configurazione obiettivo: usiamo y=7 per x=1,2,3,4,5,7 e mettiamo M in (5,7) e A in (7,7)? O forse usiamo la riga 6? Le celle valide in y=6: x=1,2,3,4,5,7 (x=6 è bordo). Quindi possiamo usare:\\ \tt P(1,6), T(2,6), S(3,6), D(4,6), A(5,6), M(7,6)? Ma O è già in (7,6). Quindi forse spostare O altrove.\\ \tt Alternative: usare riga 5? Le celle valide in y=5: x=1,2,3,4,5 (x=6 è bordo). Non abbastanza spazio.\\ \tt Meglio usare riga 7 con x=1,2,3,4,5,7: quindi P(1,7), T(2,7), S(3,7), D(4,7), A(5,7), O(7,7) e M in (6,6) o altro? Ma (6,6) è bordo.\\ \tt Propongo: P(1,7), T(2,7), S(3,7), D(4,7), A(5,7), O(7,7) e M in (7,5) o (1,6)? Dobbiamo trovare una cella libera per M.\\ \tt Sulla mia griglia, (7,5) è libero? y=5: "║◌◌◌◌◌┌─╢" so x=1-5 sono ◌, x=6 è bordo, x=7? Non è visibile, ma probabilmente è bordo. In realtà, la griglia ha x=1-7, ma i bordi occupano alcune celle.\\ \tt Per evitare confusione, propongo di usare la riga 6 per alcuni oggetti: ad esempio, mettere M in (1,6) e O in (7,6) (già lì), e gli altri in riga 7.\\ \tt Quindi obiettivo modificato:\\ \tt P(1,7), T(2,7), S(3,7), D(4,7), A(5,7), M(1,6), O(7,6). Confermi? 
	  } 
	   } 
	   } 
	 & & \\ 
 

    \theutterance \stepcounter{utterance}  

    & & & \multicolumn{2}{p{0.3\linewidth}}{\cellcolor[rgb]{0.95,0.95,0.95}{%
	\makecell[{{p{\linewidth}}}]{% 
	  \tt {\tiny [GM$|$GM]}  
	 Invalid format: tail is not empty: '\\ \tt P(1,6), T(2,6), S(3,6), D(4,6), A(5,6), M(7,6)? Ma O è già in (7,6). Quindi forse spostare O altrove.\\ \tt Alternative: usare riga 5? Le celle valide in y=5: x=1,2,3,4,5 (x=6 è bordo). Non abbastanza spazio.\\ \tt Meglio usare riga 7 con x=1,2,3,4,5,7: quindi P(1,7), T(2,7), S(3,7), D(4,7), A(5,7), O(7,7) e M in (6,6) o altro? Ma (6,6) è bordo.\\ \tt Propongo: P(1,7), T(2,7), S(3,7), D(4,7), A(5,7), O(7,7) e M in (7,5) o (1,6)? Dobbiamo trovare una cella libera per M.\\ \tt Sulla mia griglia, (7,5) è libero? y=5: "║◌◌◌◌◌┌─╢" so x=1-5 sono ◌, x=6 è bordo, x=7? Non è visibile, ma probabilmente è bordo. In realtà, la griglia ha x=1-7, ma i bordi occupano alcune celle.\\ \tt Per evitare confusione, propongo di usare la riga 6 per alcuni oggetti: ad esempio, mettere M in (1,6) e O in (7,6) (già lì), e gli altri in riga 7.\\ \tt Quindi obiettivo modificato:\\ \tt P(1,7), T(2,7), S(3,7), D(4,7), A(5,7), M(1,6), O(7,6). Confermi? 
	  } 
	   } 
	   } 
	 & & \\ 
 

    \theutterance \stepcounter{utterance}  

    & & \multicolumn{4}{p{0.6\linewidth}}{\cellcolor[rgb]{0.9,0.9,0.9}{%
	\makecell[{{p{\linewidth}}}]{% 
	  \tt {\tiny [A$\langle$GM]}  
	 Penalità: Il tuo messaggio non deve contenere nulla dopo il comando!\\ \tt \\ \tt Assicurati che la tua risposta contenga solo `COMUNICA: <MESSAGE>` o `SPOSTA: <OBJECT>, (<X>, <Y>)`, e nient'altro!\\ \tt \\ \tt You have collectively accumulated 8 of 16 penalties.\\ \tt Riprova! 
	  } 
	   } 
	   } 
	 & \\ 
 

    \theutterance \stepcounter{utterance}  

    & \multicolumn{4}{p{0.6\linewidth}}{\cellcolor[rgb]{0.8,1,0.9}{%
	\makecell[{{p{\linewidth}}}]{% 
	  \tt {\tiny [A$\rangle$GM]}  
	 COMUNICA: Propongo obiettivo modificato: P(1,7), T(2,7), S(3,7), D(4,7), A(5,7), M(1,6), O(7,6). Confermi? 
	  } 
	   } 
	   } 
	 & & \\ 
 

    \theutterance \stepcounter{utterance}  

    & & & \multicolumn{2}{p{0.3\linewidth}}{\cellcolor[rgb]{0.95,0.95,0.95}{%
	\makecell[{{p{\linewidth}}}]{% 
	  \tt {\tiny [GM$|$GM]}  
	 COMUNICA: Propongo obiettivo modificato: P(1,7), T(2,7), S(3,7), D(4,7), A(5,7), M(1,6), O(7,6). Confermi? 
	  } 
	   } 
	   } 
	 & & \\ 
 

    \theutterance \stepcounter{utterance}  

    & & \multicolumn{4}{p{0.6\linewidth}}{\cellcolor[rgb]{0.9,0.9,0.9}{%
	\makecell[{{p{\linewidth}}}]{% 
	  \tt {\tiny [GM$\rangle$B]}  
	 Il tuo messaggio è stato inoltrato all'altro giocatore.\\ \tt \\ \tt \\ \tt You are currently playing round 9 of maximum 28.\\ \tt You have collectively accumulated 8 of 16 penalties.\\ \tt L'altro giocatore ha inviato questo messaggio:\\ \tt \\ \tt \\ \tt "Propongo obiettivo modificato: P(1,7), T(2,7), S(3,7), D(4,7), A(5,7), M(1,6), O(7,6). Confermi?"\\ \tt \\ \tt \\ \tt Qual è il tuo prossimo comando? 
	  } 
	   } 
	   } 
	 & \\ 
 

    \theutterance \stepcounter{utterance}  

    & & & \multicolumn{4}{p{0.6\linewidth}}{\cellcolor[rgb]{1,0.85,0.72}{%
	\makecell[{{p{\linewidth}}}]{% 
	  \tt {\tiny [GM$\langle$B]}  
	 COMUNICA: Confermo obiettivo modificato. Sposterò M in (1,6) e O in (7,6). 
	  } 
	   } 
	   } 
	  \\ 
 

    \theutterance \stepcounter{utterance}  

    & & & \multicolumn{2}{p{0.3\linewidth}}{\cellcolor[rgb]{0.95,0.95,0.95}{%
	\makecell[{{p{\linewidth}}}]{% 
	  \tt {\tiny [GM$|$GM]}  
	 COMUNICA: Confermo obiettivo modificato. Sposterò M in (1,6) e O in (7,6). 
	  } 
	   } 
	   } 
	 & & \\ 
 

    \theutterance \stepcounter{utterance}  

    & & \multicolumn{4}{p{0.6\linewidth}}{\cellcolor[rgb]{0.9,0.9,0.9}{%
	\makecell[{{p{\linewidth}}}]{% 
	  \tt {\tiny [A$\langle$GM]}  
	 Il tuo messaggio è stato inoltrato all'altro giocatore.\\ \tt \\ \tt \\ \tt You are currently playing round 9 of maximum 28.\\ \tt You have collectively accumulated 8 of 16 penalties.\\ \tt L'altro giocatore ha inviato questo messaggio:\\ \tt \\ \tt \\ \tt "Confermo obiettivo modificato. Sposterò M in (1,6) e O in (7,6)."\\ \tt \\ \tt \\ \tt Qual è il tuo prossimo comando? 
	  } 
	   } 
	   } 
	 & \\ 
 

    \theutterance \stepcounter{utterance}  

    & \multicolumn{4}{p{0.6\linewidth}}{\cellcolor[rgb]{0.8,1,0.9}{%
	\makecell[{{p{\linewidth}}}]{% 
	  \tt {\tiny [A$\rangle$GM]}  
	 SPOSTA: A, (5, 7) 
	  } 
	   } 
	   } 
	 & & \\ 
 

    \theutterance \stepcounter{utterance}  

    & & & \multicolumn{2}{p{0.3\linewidth}}{\cellcolor[rgb]{0.95,0.95,0.95}{%
	\makecell[{{p{\linewidth}}}]{% 
	  \tt {\tiny [GM$|$GM]}  
	 SPOSTA: A, (5, 7) 
	  } 
	   } 
	   } 
	 & & \\ 
 

    \theutterance \stepcounter{utterance}  

    & & & \multicolumn{2}{p{0.3\linewidth}}{\cellcolor[rgb]{0.95,0.95,0.95}{%
	\makecell[{{p{\linewidth}}}]{% 
	  \tt {\tiny [GM$|$GM]}  
	 Spostato 'A' in (5,7) con successo. La tua griglia aggiornata appare così:\\ \tt ```\\ \tt  1234567\\ \tt ╔═╤═╤═╤═╗\\ \tt ║◌│◌│◌│◌║ 1\\ \tt ╟─┴─┼─┤◌║ 2\\ \tt ║◌◌M│◌└─╢ 3\\ \tt ╟───┘◌◌◌║ 4\\ \tt ║◌◌◌◌◌┌─╢ 5\\ \tt ║◌◌◌◌◌│O║ 6\\ \tt ║PTSDA│◌║ 7\\ \tt ╚═════╧═╝\\ \tt ```\\ \tt  
	  } 
	   } 
	   } 
	 & & \\ 
 

    \theutterance \stepcounter{utterance}  

    & & \multicolumn{4}{p{0.6\linewidth}}{\cellcolor[rgb]{0.9,0.9,0.9}{%
	\makecell[{{p{\linewidth}}}]{% 
	  \tt {\tiny [GM$\rangle$B]}  
	 Il tuo messaggio è stato inoltrato all'altro giocatore.\\ \tt \\ \tt \\ \tt You are currently playing round 10 of maximum 28.\\ \tt You have collectively accumulated 8 of 16 penalties.\\ \tt L'altro giocatore ha spostato un oggetto sulla sua griglia.\\ \tt \\ \tt \\ \tt Qual è il tuo prossimo comando? 
	  } 
	   } 
	   } 
	 & \\ 
 

    \theutterance \stepcounter{utterance}  

    & & & \multicolumn{4}{p{0.6\linewidth}}{\cellcolor[rgb]{1,0.85,0.72}{%
	\makecell[{{p{\linewidth}}}]{% 
	  \tt {\tiny [GM$\langle$B]}  
	 SPOSTA: M, (1,6) 
	  } 
	   } 
	   } 
	  \\ 
 

    \theutterance \stepcounter{utterance}  

    & & & \multicolumn{2}{p{0.3\linewidth}}{\cellcolor[rgb]{0.95,0.95,0.95}{%
	\makecell[{{p{\linewidth}}}]{% 
	  \tt {\tiny [GM$|$GM]}  
	 SPOSTA: M, (1,6) 
	  } 
	   } 
	   } 
	 & & \\ 
 

    \theutterance \stepcounter{utterance}  

    & & & \multicolumn{2}{p{0.3\linewidth}}{\cellcolor[rgb]{0.95,0.95,0.95}{%
	\makecell[{{p{\linewidth}}}]{% 
	  \tt {\tiny [GM$|$GM]}  
	 Spostato 'M' in (1,6) con successo. La tua griglia aggiornata appare così:\\ \tt ```\\ \tt  1234567\\ \tt ╔═╤═╤═╤═╗\\ \tt ║◌│◌│◌│◌║ 1\\ \tt ╟─┴─┼─┤◌║ 2\\ \tt ║◌◌◌│◌└─╢ 3\\ \tt ╟───┘◌◌◌║ 4\\ \tt ║◌◌◌◌◌┌─╢ 5\\ \tt ║M◌◌◌◌│◌║ 6\\ \tt ║PTSDA│O║ 7\\ \tt ╚═════╧═╝\\ \tt ```\\ \tt  
	  } 
	   } 
	   } 
	 & & \\ 
 

    \theutterance \stepcounter{utterance}  

    & & \multicolumn{4}{p{0.6\linewidth}}{\cellcolor[rgb]{0.9,0.9,0.9}{%
	\makecell[{{p{\linewidth}}}]{% 
	  \tt {\tiny [A$\langle$GM]}  
	 Spostato 'A' in (5,7) con successo. La tua griglia aggiornata appare così:\\ \tt ```\\ \tt  1234567\\ \tt ╔═╤═╤═╤═╗\\ \tt ║◌│◌│◌│◌║ 1\\ \tt ╟─┴─┼─┤◌║ 2\\ \tt ║◌◌M│◌└─╢ 3\\ \tt ╟───┘◌◌◌║ 4\\ \tt ║◌◌◌◌◌┌─╢ 5\\ \tt ║◌◌◌◌◌│O║ 6\\ \tt ║PTSDA│◌║ 7\\ \tt ╚═════╧═╝\\ \tt ```\\ \tt You are currently playing round 10 of maximum 28.\\ \tt You have collectively accumulated 8 of 16 penalties.\\ \tt L'altro giocatore ha spostato un oggetto sulla sua griglia.\\ \tt \\ \tt \\ \tt Qual è il tuo prossimo comando? 
	  } 
	   } 
	   } 
	 & \\ 
 

    \theutterance \stepcounter{utterance}  

    & \multicolumn{4}{p{0.6\linewidth}}{\cellcolor[rgb]{0.8,1,0.9}{%
	\makecell[{{p{\linewidth}}}]{% 
	  \tt {\tiny [A$\rangle$GM]}  
	 SPOSTA: M, (1, 6) 
	  } 
	   } 
	   } 
	 & & \\ 
 

    \theutterance \stepcounter{utterance}  

    & & & \multicolumn{2}{p{0.3\linewidth}}{\cellcolor[rgb]{0.95,0.95,0.95}{%
	\makecell[{{p{\linewidth}}}]{% 
	  \tt {\tiny [GM$|$GM]}  
	 SPOSTA: M, (1, 6) 
	  } 
	   } 
	   } 
	 & & \\ 
 

    \theutterance \stepcounter{utterance}  

    & & & \multicolumn{2}{p{0.3\linewidth}}{\cellcolor[rgb]{0.95,0.95,0.95}{%
	\makecell[{{p{\linewidth}}}]{% 
	  \tt {\tiny [GM$|$GM]}  
	 Spostato 'M' in (1,6) con successo. La tua griglia aggiornata appare così:\\ \tt ```\\ \tt  1234567\\ \tt ╔═╤═╤═╤═╗\\ \tt ║◌│◌│◌│◌║ 1\\ \tt ╟─┴─┼─┤◌║ 2\\ \tt ║◌◌◌│◌└─╢ 3\\ \tt ╟───┘◌◌◌║ 4\\ \tt ║◌◌◌◌◌┌─╢ 5\\ \tt ║M◌◌◌◌│O║ 6\\ \tt ║PTSDA│◌║ 7\\ \tt ╚═════╧═╝\\ \tt ```\\ \tt  
	  } 
	   } 
	   } 
	 & & \\ 
 

    \theutterance \stepcounter{utterance}  

    & & \multicolumn{4}{p{0.6\linewidth}}{\cellcolor[rgb]{0.9,0.9,0.9}{%
	\makecell[{{p{\linewidth}}}]{% 
	  \tt {\tiny [GM$\rangle$B]}  
	 Spostato 'M' in (1,6) con successo. La tua griglia aggiornata appare così:\\ \tt ```\\ \tt  1234567\\ \tt ╔═╤═╤═╤═╗\\ \tt ║◌│◌│◌│◌║ 1\\ \tt ╟─┴─┼─┤◌║ 2\\ \tt ║◌◌◌│◌└─╢ 3\\ \tt ╟───┘◌◌◌║ 4\\ \tt ║◌◌◌◌◌┌─╢ 5\\ \tt ║M◌◌◌◌│◌║ 6\\ \tt ║PTSDA│O║ 7\\ \tt ╚═════╧═╝\\ \tt ```\\ \tt You are currently playing round 11 of maximum 28.\\ \tt You have collectively accumulated 8 of 16 penalties.\\ \tt L'altro giocatore ha spostato un oggetto sulla sua griglia.\\ \tt \\ \tt \\ \tt Qual è il tuo prossimo comando? 
	  } 
	   } 
	   } 
	 & \\ 
 

    \theutterance \stepcounter{utterance}  

    & & & \multicolumn{4}{p{0.6\linewidth}}{\cellcolor[rgb]{1,0.85,0.72}{%
	\makecell[{{p{\linewidth}}}]{% 
	  \tt {\tiny [GM$\langle$B]}  
	 SPOSTA: O, (7,6) 
	  } 
	   } 
	   } 
	  \\ 
 

    \theutterance \stepcounter{utterance}  

    & & & \multicolumn{2}{p{0.3\linewidth}}{\cellcolor[rgb]{0.95,0.95,0.95}{%
	\makecell[{{p{\linewidth}}}]{% 
	  \tt {\tiny [GM$|$GM]}  
	 SPOSTA: O, (7,6) 
	  } 
	   } 
	   } 
	 & & \\ 
 

    \theutterance \stepcounter{utterance}  

    & & & \multicolumn{2}{p{0.3\linewidth}}{\cellcolor[rgb]{0.95,0.95,0.95}{%
	\makecell[{{p{\linewidth}}}]{% 
	  \tt {\tiny [GM$|$GM]}  
	 Spostato 'O' in (7,6) con successo. La tua griglia aggiornata appare così:\\ \tt ```\\ \tt  1234567\\ \tt ╔═╤═╤═╤═╗\\ \tt ║◌│◌│◌│◌║ 1\\ \tt ╟─┴─┼─┤◌║ 2\\ \tt ║◌◌◌│◌└─╢ 3\\ \tt ╟───┘◌◌◌║ 4\\ \tt ║◌◌◌◌◌┌─╢ 5\\ \tt ║M◌◌◌◌│O║ 6\\ \tt ║PTSDA│◌║ 7\\ \tt ╚═════╧═╝\\ \tt ```\\ \tt  
	  } 
	   } 
	   } 
	 & & \\ 
 

    \theutterance \stepcounter{utterance}  

    & & \multicolumn{4}{p{0.6\linewidth}}{\cellcolor[rgb]{0.9,0.9,0.9}{%
	\makecell[{{p{\linewidth}}}]{% 
	  \tt {\tiny [A$\langle$GM]}  
	 Spostato 'M' in (1,6) con successo. La tua griglia aggiornata appare così:\\ \tt ```\\ \tt  1234567\\ \tt ╔═╤═╤═╤═╗\\ \tt ║◌│◌│◌│◌║ 1\\ \tt ╟─┴─┼─┤◌║ 2\\ \tt ║◌◌◌│◌└─╢ 3\\ \tt ╟───┘◌◌◌║ 4\\ \tt ║◌◌◌◌◌┌─╢ 5\\ \tt ║M◌◌◌◌│O║ 6\\ \tt ║PTSDA│◌║ 7\\ \tt ╚═════╧═╝\\ \tt ```\\ \tt You are currently playing round 11 of maximum 28.\\ \tt You have collectively accumulated 8 of 16 penalties.\\ \tt L'altro giocatore ha spostato un oggetto sulla sua griglia.\\ \tt \\ \tt \\ \tt Qual è il tuo prossimo comando? 
	  } 
	   } 
	   } 
	 & \\ 
 

    \theutterance \stepcounter{utterance}  

    & \multicolumn{4}{p{0.6\linewidth}}{\cellcolor[rgb]{0.8,1,0.9}{%
	\makecell[{{p{\linewidth}}}]{% 
	  \tt {\tiny [A$\rangle$GM]}  
	 COMUNICA: Ho spostato A(5,7) e M(1,6). O è già in (7,6). Tutti gli oggetti sono nella posizione obiettivo? Controllo: P(1,7), T(2,7), S(3,7), D(4,7), A(5,7), M(1,6), O(7,6). Sulla mia griglia, è così. Sei pronto per finire? 
	  } 
	   } 
	   } 
	 & & \\ 
 

    \theutterance \stepcounter{utterance}  

    & & & \multicolumn{2}{p{0.3\linewidth}}{\cellcolor[rgb]{0.95,0.95,0.95}{%
	\makecell[{{p{\linewidth}}}]{% 
	  \tt {\tiny [GM$|$GM]}  
	 COMUNICA: Ho spostato A(5,7) e M(1,6). O è già in (7,6). Tutti gli oggetti sono nella posizione obiettivo? Controllo: P(1,7), T(2,7), S(3,7), D(4,7), A(5,7), M(1,6), O(7,6). Sulla mia griglia, è così. Sei pronto per finire? 
	  } 
	   } 
	   } 
	 & & \\ 
 

    \theutterance \stepcounter{utterance}  

    & & \multicolumn{4}{p{0.6\linewidth}}{\cellcolor[rgb]{0.9,0.9,0.9}{%
	\makecell[{{p{\linewidth}}}]{% 
	  \tt {\tiny [GM$\rangle$B]}  
	 Spostato 'O' in (7,6) con successo. La tua griglia aggiornata appare così:\\ \tt ```\\ \tt  1234567\\ \tt ╔═╤═╤═╤═╗\\ \tt ║◌│◌│◌│◌║ 1\\ \tt ╟─┴─┼─┤◌║ 2\\ \tt ║◌◌◌│◌└─╢ 3\\ \tt ╟───┘◌◌◌║ 4\\ \tt ║◌◌◌◌◌┌─╢ 5\\ \tt ║M◌◌◌◌│O║ 6\\ \tt ║PTSDA│◌║ 7\\ \tt ╚═════╧═╝\\ \tt ```\\ \tt You are currently playing round 12 of maximum 28.\\ \tt You have collectively accumulated 8 of 16 penalties.\\ \tt L'altro giocatore ha inviato questo messaggio:\\ \tt \\ \tt \\ \tt "Ho spostato A(5,7) e M(1,6). O è già in (7,6). Tutti gli oggetti sono nella posizione obiettivo? Controllo: P(1,7), T(2,7), S(3,7), D(4,7), A(5,7), M(1,6), O(7,6). Sulla mia griglia, è così. Sei pronto per finire?"\\ \tt \\ \tt \\ \tt Qual è il tuo prossimo comando? 
	  } 
	   } 
	   } 
	 & \\ 
 

    \theutterance \stepcounter{utterance}  

    & & & \multicolumn{4}{p{0.6\linewidth}}{\cellcolor[rgb]{1,0.85,0.72}{%
	\makecell[{{p{\linewidth}}}]{% 
	  \tt {\tiny [GM$\langle$B]}  
	 COMUNICA: finito? 
	  } 
	   } 
	   } 
	  \\ 
 

    \theutterance \stepcounter{utterance}  

    & & & \multicolumn{2}{p{0.3\linewidth}}{\cellcolor[rgb]{0.95,0.95,0.95}{%
	\makecell[{{p{\linewidth}}}]{% 
	  \tt {\tiny [GM$|$GM]}  
	 COMUNICA: finito? 
	  } 
	   } 
	   } 
	 & & \\ 
 

    \theutterance \stepcounter{utterance}  

    & & \multicolumn{4}{p{0.6\linewidth}}{\cellcolor[rgb]{0.9,0.9,0.9}{%
	\makecell[{{p{\linewidth}}}]{% 
	  \tt {\tiny [A$\langle$GM]}  
	 Il tuo messaggio è stato inoltrato all'altro giocatore.\\ \tt \\ \tt \\ \tt You are currently playing round 12 of maximum 28.\\ \tt You have collectively accumulated 8 of 16 penalties.\\ \tt L'altro giocatore ha inviato questo messaggio:\\ \tt \\ \tt \\ \tt "finito?"\\ \tt \\ \tt \\ \tt Qual è il tuo prossimo comando? 
	  } 
	   } 
	   } 
	 & \\ 
 

    \theutterance \stepcounter{utterance}  

    & \multicolumn{4}{p{0.6\linewidth}}{\cellcolor[rgb]{0.8,1,0.9}{%
	\makecell[{{p{\linewidth}}}]{% 
	  \tt {\tiny [A$\rangle$GM]}  
	 COMUNICA: finito! 
	  } 
	   } 
	   } 
	 & & \\ 
 

    \theutterance \stepcounter{utterance}  

    & & & \multicolumn{2}{p{0.3\linewidth}}{\cellcolor[rgb]{0.95,0.95,0.95}{%
	\makecell[{{p{\linewidth}}}]{% 
	  \tt {\tiny [GM$|$GM]}  
	 COMUNICA: finito! 
	  } 
	   } 
	   } 
	 & & \\ 
 

    \theutterance \stepcounter{utterance}  

    & & & \multicolumn{2}{p{0.3\linewidth}}{\cellcolor[rgb]{0.95,0.95,0.95}{%
	\makecell[{{p{\linewidth}}}]{% 
	  \tt {\tiny [GM$|$GM]}  
	 true 
	  } 
	   } 
	   } 
	 & & \\ 
 

    \theutterance \stepcounter{utterance}  

    & & & \multicolumn{2}{p{0.3\linewidth}}{\cellcolor[rgb]{0.95,0.95,0.95}{%
	\makecell[{{p{\linewidth}}}]{% 
	  \tt {\tiny [GM$|$GM]}  
	 * success: True\\ \tt * lose: False\\ \tt * aborted: False\\ \tt -------\\ \tt * Shifts: 10.00\\ \tt * Max Shifts: 12.00\\ \tt * Min Shifts: 6.00\\ \tt * End Distance Sum: 0.00\\ \tt * Init Distance Sum: 18.94\\ \tt * Expected Distance Sum: 29.33\\ \tt * Penalties: 8.00\\ \tt * Max Penalties: 16.00\\ \tt * Rounds: 12.00\\ \tt * Max Rounds: 28.00\\ \tt * Object Count: 7.00\\ \tt  
	  } 
	   } 
	   } 
	 & & \\ 
 

    \theutterance \stepcounter{utterance}  

    & & & \multicolumn{2}{p{0.3\linewidth}}{\cellcolor[rgb]{0.95,0.95,0.95}{%
	\makecell[{{p{\linewidth}}}]{% 
	  \tt {\tiny [GM$|$GM]}  
	 * Main Score: 40.21\\ \tt \\ \tt -------\\ \tt * Distance Score: 1.00\\ \tt * Consistency Score: 0.43\\ \tt * Coverage Score: 0.88\\ \tt * Penalty Score: 0.60\\ \tt * Alternative Penalty Score: 0.75\\ \tt * Alternative Main Score: 75.00\\ \tt \\ \tt -------\\ \tt * Shifts: 10.00\\ \tt * Max Shifts: 12.00\\ \tt * Min Shifts: 6.00\\ \tt * End Distance Sum: 0.00\\ \tt * Init Distance Sum: 18.94\\ \tt * Expected Distance Sum: 29.33\\ \tt * Penalties: 8.00\\ \tt * Max Penalties: 16.00\\ \tt * Rounds: 12.00\\ \tt * Max Rounds: 28.00\\ \tt * Object Count: 7.00\\ \tt  
	  } 
	   } 
	   } 
	 & & \\ 
 

\end{supertabular}
}

\end{document}
